\documentclass[11pt]{article}

% --- 1. 基础宏包 ---
\usepackage[margin=1.6cm]{geometry}
\usepackage{amsmath,amssymb,amsthm}
\allowdisplaybreaks[4]
\usepackage[hidelinks]{hyperref}
\hypersetup{
	colorlinks=true,      % 激活彩色链接
	linkcolor=blue,       % 内部链接（如目录、公式）的颜色
	filecolor=blue,       % 本地文件的链接颜色
	urlcolor=blue,        % 网页 URL 的颜色
	citecolor=blue,       % 参考文献引用的颜色
}
\usepackage[most]{tcolorbox}
\usepackage{cancel}
\usepackage{tikz-feynman}
\tikzfeynmanset{compat=1.1.0}
\usepackage{slashed}
\usepackage{simpler-wick}
\usepackage{braket}

% --- 2. 整合专用设置 ---
\usepackage{docmute}   % 自动忽略子文档的导言区
\usepackage{tocloft}   
\usepackage{totcount}  
\usepackage{xparse}    

% --- 3. 目录与样式设置 ---
\newcommand{\listproblemname}{List of Solved Problems}
\newlistof{myproblem}{los}{\listproblemname}
\setlength{\cftbeforemyproblemskip}{5pt}

\newcounter{numSolved}
\regtotcounter{numSolved}
\newcommand{\totalQuestions}{80}

% --- 4. Solution 环境定义 (主文档逻辑) ---
% 当整合编译时，使用这个带有计数和目录功能的定义
\NewDocumentEnvironment{solution}{o}
{
	\begin{proof}[\textbf{Solution}]
		\IfValueT{#1}{%
			\stepcounter{numSolved}%
			\phantomsection 
			\addcontentsline{los}{myproblem}{Problem #1}%
			\textbf{(#1) }
		}
	}
	{
	\end{proof}
}

% 保持子文档的题目框样式
\tcbset{
	problem/.style={
		enhanced,
		boxrule=0.7pt,
		arc=3mm,
		left=3mm,right=3mm,top=2mm,bottom=2mm,
		colback=white,
		colframe=black,
		before skip=10pt,
		after skip=6pt,
	}
}

\title{\Huge Solution Manual of QFT Assignments\\(Problem Set 11--18)}
\author{Bufan Zheng\\ \href{mailto:bufan@g.ecc.u-tokyo.ac.jp}{\texttt{bufan@g.ecc.u-tokyo.ac.jp}} }
\date{}

\begin{document}
	
	\maketitle
	
	% 进度展示
	\begin{center}
		\Large\textbf{Total Progress: \total{numSolved} / \totalQuestions}
	\end{center}
	
	% 目录
	\listofmyproblem
	
	% --- 5. 导入子文档 ---
	% 使用 clearpage 确保每个文档从新页面开始
	\clearpage
	\documentclass[11pt]{article}

\usepackage[margin=1.6cm]{geometry}
\usepackage{amsmath,amssymb,amsthm}
\usepackage[hidelinks]{hyperref}
\usepackage[most]{tcolorbox}
\usepackage[]{cancel}
\usepackage{tikz-feynman}
\tikzfeynmanset{compat=1.1.0}
% --- Rounded box style for each problem (whole statement) ---
\tcbset{
	problem/.style={
		enhanced,
		boxrule=0.7pt,
		arc=3mm,
		left=3mm,right=3mm,top=2mm,bottom=2mm,
		colback=white,
		colframe=black,
		before skip=10pt,
		after skip=6pt,
	}
}

% --- "Solution" environment (proof-style) ---
\newenvironment{solution}[1][]{\begin{proof}[\textbf{Solution}]}{\end{proof}}


\begin{document}
\begin{center}
    {\LARGE \textbf{Solution Manual of Problem Set 11} \par}
    \vspace{1em}
    {\large Bufan Zheng\\ \href{mailto:bufan@g.ecc.u-tokyo.ac.jp}{\texttt{bufan@g.ecc.u-tokyo.ac.jp}} \par}
\end{center}

	\begin{tcolorbox}[problem]
		\textbf{1.} Starting from \(Z_0[J] = \exp\!\left(\frac{i}{2}\int J\,\Delta_F\,J\right)\), show that
		\(\langle 0|T\,\phi(x)\phi(y)|0\rangle_0=\Delta_F(x-y)\) by functional differentiation.
	\end{tcolorbox}
	\begin{solution}[11.1]
		\[
		\begin{aligned}
			\langle 0|T\phi(x)\phi(y)|0\rangle_0
			&=\left.\frac{1}{Z_0[J]}\frac{\delta}{i\delta J(x)}\frac{\delta}{i\delta J(y)}\,Z_0[J]\right|_{J=0} \\
			&=\left.\frac{1}{Z_0[J]}\frac{\delta}{i\delta J(x)}\frac{\delta}{i\delta J(y)}
			\exp\!\left[\frac{i}{2}\int d^4z\,d^4w\,J(z)\,\Delta_F(z-w)\,J(w)\right]\right|_{J=0} \\
			&=\left.\frac{1}{Z_0[J]}\frac{\delta}{i\delta J(x)}
			\left(\int d^4z\,\Delta_F(z-y)\,J(z)\,Z_0[J]\right)\right|_{J=0} \\
			&=\left.-i\,\frac{1}{\cancel{Z_0[J]}}\,\Delta_F(x-y)\,\cancel{Z_0[J]}\right|_{J=0}
			+\left.\frac{1}{Z_0[J]}
			\cancelto{0}{\left(\int d^4z\,\Delta_F(z-y)\,J(z)\right)^{2}} Z_0[J]\right|_{J=0} \\
			&=-i\,\Delta_F(x-y).
		\end{aligned}
		\]
		So \textbf{MY FINAL ANSWER IS}:
		\[
		\boxed{
		i\langle 0|T\,\phi(x)\phi(y)|0\rangle_0=\Delta_F(x-y)
		}
		\]
	\end{solution}
	
	\begin{tcolorbox}[problem]
		\textbf{2.} In \(\phi^4\) theory with \(\mathcal L_{\text{int}}=-\frac{\lambda}{4!}\phi^4\), expand \(Z[J]\) to first order in \(\lambda\) and identify the term that produces the ``figure-8'' vacuum bubble.
	\end{tcolorbox}
	\begin{solution}[11.2]
		\[
		\begin{aligned}
			Z[J]=&\int\mathcal{D}\phi e^{i\int d^d x\left[\mathcal{L}_0+\mathcal{L}\_1+J\phi\right]}=e^{i\int d^d x\mathcal{L}_1\left(\frac{\delta}{i\delta J(x)}\right)}\int \mathcal{D}\phi e^{i\int d^d x\left[\mathcal{L}_0+J\phi\right]}\\
			=&\large e^{\large i\int d^dx\left.\mathcal{L}_1\left(\frac{1}{i}\frac{\delta}{\delta J(x)}\right)\right.}Z_0(J)=\exp\left[-\frac{i\lambda}{4!}\int d^dx\left(\frac{1}{i}\frac{\delta}{\delta J(x)}\right)^4\right]Z_0(J)\\
			=&Z_0[J]-\left[\frac{i\lambda}{4!}\int d^dx\left(\frac{1}{i}\frac{\delta}{\delta J(x)}\right)^4\right]Z_0(J)+\mathcal{O}(\lambda^2)
		\end{aligned}
		\]
	\end{solution}
	from last problem we know that:
	\[
	\frac 1i\frac{\delta}{\delta J(x)}Z_0[J]=\int d^d z\Delta_F(z-x)J(z)Z_0[J]
	\]
	Therefore, we have:
	% (若要显示被划掉的那一项，请在导言区加 \usepackage{cancel})
	\[
	\begin{aligned}
		\left(\frac{\delta}{i\,\delta J(x)}\right)^{4} Z_0[J]
		&=\left(\frac{\delta}{i\,\delta J(x)}\right)^{3}\int d^{4}z_{1}\,\Delta_{F}(z_{1}-x)\,J(z_{1})\,Z_{0}[J] \\
		&=\left(\frac{\delta}{i\,\delta J(x)}\right)^{2}\Bigl(\Delta_{F}(0)\,Z_{0}[J]
		+\int d^{4}z_{1}\,d^{4}z_{2}\,\Delta_{F}(z_{1}-x)\Delta_{F}(z_{2}-x)\,J(z_{1})J(z_{2})\,Z_{0}[J]\Bigr) \\
		&=\left(\frac{\delta}{i\,\delta J(x)}\right)\Bigl(\Delta_{F}(0)\int d^{4}z_{1}\,\Delta_{F}(z_{1}-x)J(z_1)\,Z_{0}[J]
		+2\int d^{4}z_{1}\,\Delta_{F}(z_{1}-x)\,\Delta_{F}(0)\,J(z_{1})\,Z_{0}[J] \\
		&\hspace{3.6em}
		+\int d^{4}z_{1}\,d^{4}z_{2}\,d^{4}z_{3}\,\Delta_{F}(z_{1}-x)\Delta_{F}(z_{2}-x)\Delta_{F}(z_{3}-x)\,
		J(z_{1})J(z_{2})J(z_{3})\,Z_{0}[J]\Bigr) \\
		&=\Delta_{F}^{2}(0)\,Z_{0}[J]
		+\Delta_{F}(0)\int d^{4}z_{1}\,d^{4}z_{2}\,\Delta_{F}(z_{1}-x)\Delta_{F}(z_{2}-x)\,J(z_{1})J(z_{2})\,Z_{0}[J] \\
		&\quad
		+2\int d^{4}z_{1}\,d^{4}z_{2}\,\Delta_{F}(z_{1}-x)\Delta_{F}(z_{2}-x)\,\Delta_{F}(0)\,J(z_{1})J(z_{2})\,Z_{0}[J] \\
		&\quad
		+2\Delta_{F}^{2}(0)\,Z_{0}[J]		+3\int d^{4}z_{1}\,d^{4}z_{2}\,\Delta_{F}(z_{1}-x)\Delta_{F}(z_{2}-x)\,J(z_{1})J(z_{2})\,\Delta_{F}(0)\,Z_{0}[J]  \\
		&\quad
		+\int d^{4}z_{1}\,d^{4}z_{2}\,d^{4}z_{3}\,d^{4}z_{4}\,
		\Delta_{F}(z_{1}-x)\Delta_{F}(z_{2}-x)\Delta_{F}(z_{3}-x)\Delta_{F}(z_{4}-x)\,
		J(z_{1})J(z_{2})J(z_{3})J(z_{4})\,Z_{0}[J]\\
		&=Z_{0}[J]\Bigl(3\Delta_{F}^{2}(0)
		+6\Delta_{F}(0)\int d^{4}z_{1}\,d^{4}z_{2}\,\Delta_{F}(z_{1}-x)\Delta_{F}(z_{2}-x)\,J(z_{1})J(z_{2}) \\
		&\quad
		+\int d^{4}z_{1}\,d^{4}z_{2}\,d^{4}z_{3}\,d^{4}z_{4}\,
		\Delta_{F}(z_{1}-x)\Delta_{F}(z_{2}-x)\Delta_{F}(z_{3}-x)\Delta_{F}(z_{4}-x)\,
		J(z_{1})J(z_{2})J(z_{3})J(z_{4})\Bigr).
	\end{aligned}
	\]
	For the vacuume diagram, we can set $J=0$:
	\[
	Z[0]=1-i\frac{\lambda}{8}\int d^4x\Delta_F^2(0)+\mathcal{O}(\lambda^2)
	\]
	Hence, \textbf{MY FINAL ANSWER IS}:
\[
\boxed{
\tikz[baseline={(v.base)}, line width=0.9pt]{
	\coordinate (v) at (0,0);
	\draw (v) .. controls (-1, 1) and (-1,-1) .. (v);
	\draw (v) .. controls ( 1, 1) and ( 1,-1) .. (v);
	\fill (v) circle (1.2pt);
}
\;=\;
\frac{1}{8}\cdot\left(-i\lambda\right)\cdot\int d^4x\,\Delta_F(0)^2
}
\]

	
	\begin{tcolorbox}[problem]
		\textbf{3.} Derive the momentum-space propagator
		\(\Delta_F(p)=\dfrac{i}{p^2+m^2-i\epsilon}\)
		from the quadratic action by Fourier transform and inversion.
	\end{tcolorbox}
	\begin{solution}[11.3]
		From the Lagrangian, we know that:
		\[
		\mathcal{L_0}=\frac12\partial^\mu\phi\partial_\mu\phi+\frac12m^2\phi^2\sim\frac12\phi^2(-\partial^2+m^2)\phi^2
		\]
		It means  we need to solve the following equation to find the Green function:
		\[
		(-\partial^2+m^2-i\epsilon)\Delta_F(x-y)=i\mathrm{~}\delta^{(4)}(x-y)
		\]
		Using the Fourier transformation, we have:
		\[
		(p^2+m^2-i\epsilon)\Delta_F(p)=i\
		\]
			Hence, \textbf{MY FINAL ANSWER IS}:
		\[
		\boxed{
		\Delta_F(p)=\frac{i}{p^2+m^2-i\epsilon}
	}
		\]
	\end{solution}

	\begin{tcolorbox}[problem]
		\textbf{4.} Using the path-integral expansion, explain why the \(\phi^4\) vertex factor is \(-i\lambda\) and why a
		\((2\pi)^D\delta^{(D)}\!\left(\sum p\right)\) appears.
	\end{tcolorbox}
	\begin{solution}[11.4]
		We can use the calculations in problem 2. There is a term in $Z[J]$:
		\[
		Z[J]\supset -\frac{i\lambda}{4!}\int d^4x\int d^{4}z_{1}\,d^{4}z_{2}\,d^{4}z_{3}\,d^{4}z_{4}\,
		\Delta_{F}(z_{1}-x)\Delta_{F}(z_{2}-x)\Delta_{F}(z_{3}-x)\Delta_{F}(z_{4}-x)\,
		J(z_{1})J(z_{2})J(z_{3})J(z_{4})
		\]
		To get the Feynman rule for $\phi^4$ vertex factor, we need to  apply $\Pi_{i=1}^4 \frac{\delta}{i\delta J(z_i)}$ on it. It is clear there are $4!$ different way, which will cancel the front factor $\frac{1}{4!}$. Then, we obtain the vertex in coordinate space:
		\[
		V({z_1,z_2,z_3,z_4})= - i\lambda \int d^4x \Delta_F(z_1-x)\Delta_F(z_2-x)\Delta_F(z_3-x)\Delta_F(z_4-x)
		\]
		Next, do the Fourie transformation to work in momentum space:
		\[
		\begin{aligned}
		V({p_1,p_2,p_3,p_4})&= - i\lambda\int d^4z_1 d^4z_2 d^4z_3d^4z_4 e^{ip_1\cdot z_1}e^{ip_2\cdot z_2}e^{ip_3\cdot z_3}e^{ip_4\cdot z_4}\int d^4x \Delta_F(z_1-x)\Delta_F(z_2-x)\Delta_F(z_3-x)\Delta_F(z_4-x)\\
		&=-i\lambda\cdot \Delta_F(p_1)\Delta_F(p_2) \Delta_F(p_3) \Delta_F(p_4)  \cdot\int d^4 x e^{ix\cdot\sum p}\\
				&=-i\lambda\cdot \Delta_F(p_1)\Delta_F(p_2) \Delta_F(p_3) \Delta_F(p_4)  \cdot(2\pi)^4\delta^{4}(\sum p)
		\end{aligned}
		\]
		 You can see the $-i\lambda$ and ${(2\pi)^D\delta^{(D)}(\sum p)}$ appear. Here, I just set $D=4$. Now, we can get rid of four external legs, i.e. $\Delta_F(p_i)$. So, \textbf{MY FINAL ANSWER IS}:
		 \[
		 \boxed{
		 	V({p_1,p_2,p_3,p_4})= - i\lambda\cdot{(2\pi)^D\delta^{(D)}(\sum p)}
		 }
		 \]
	\end{solution}
	
	\begin{tcolorbox}[problem]
		\textbf{5.} Compute the symmetry factor of the one-loop correction to the 2-point function (``tadpole'') in \(\phi^4\) theory.
	\end{tcolorbox}
		\begin{solution}[11.5]
	For the one-loop correction to the 2-point function with a single \(\phi^4\) vertex, two legs of the vertex attach to the two external points \(x\) and \(y\), and the remaining two legs contract with each other to form the loop.
	
	Counting Wick contractions:
	\begin{enumerate}
		\item Choose which leg of the 4-legged vertex connects to the external point \(x\): \(4\) choices.
		\item Choose which remaining leg connects to the external point \(y\): \(3\) choices.
		\item The two leftover legs must contract with each other to make the loop: \(1\) choice.
	\end{enumerate}
	So the number of distinct contractions is \(4\times 3\). The vertex from the interaction brings the factor \(1/4!\), hence the overall combinatorial factor is
	\[
	\frac{4\times 3}{4!}=\frac{12}{24}=\frac{1}{2}.
	\]
	Therefore the symmetry factor is \(S=2\). So, \textbf{MY FINAL ANSWER IS}:
		\[
		\boxed{
		\tikz[baseline={(v.base)}, line width=0.9pt]{
			\coordinate (v) at (0,0);
			
			% external legs
			\draw (-1,0) -- (v);
			\draw (v) -- (1,0);
			
			% loop attached to the same vertex
			\draw (v) .. controls (-0.9,1.2) and (0.9,1.2) .. (v);
			
			% vertex dot
			\fill (v) circle (1.2pt);
			
			% optional labels (remove if you do not want them)
			\node[left]  at (-1,0) {$x$};
			\node[right] at ( 1,0) {$y$};
		}
		\quad \Rightarrow\quad S=\frac{1}{2}
	}
		\]
	\end{solution}
	
	\begin{tcolorbox}[problem]
		\textbf{6.} In scalar QED, expand \(|D_\mu\phi|^2\) and explicitly list the interaction terms proportional to \(e\) and \(e^2\).
		Identify which term gives the ``seagull'' vertex.
	\end{tcolorbox}
	\begin{solution}[11.6]
	\[
	|D_\mu\phi|^2
	=(\partial_\mu\phi^\ast-i e A_\mu\phi^\ast)(\partial^\mu\phi+i e A^\mu\phi)
	=\partial_\mu\phi^\ast\,\partial^\mu\phi
	+i e\,A^\mu\bigl(\phi\,\partial_\mu\phi^\ast-\phi^\ast\,\partial_\mu\phi\bigr)
	+e^2\,A_\mu A^\mu\,\phi^\ast\phi
	\]
	\textbf{MY FINAL ANSWER IS}:
	\begin{itemize}
		\item $\propto e$, $\mathcal L_{\text{int}}^{(e)}
		= i e\,A^\mu\bigl(\phi\,\partial_\mu\phi^\ast-\phi^\ast\,\partial_\mu\phi\bigr)$:
		\[
		\feynmandiagram [baseline=(v.base), horizontal=a to b] {
			a [particle=\(\phi(p)\)] -- [charged scalar,momentum=\(p\)] v -- [charged scalar,momentum=\(p'\)] b [particle=\(\phi(p')\)],
			v -- [photon] c [particle=\(A_\mu(q)\)],
		};
		\qquad\Rightarrow\qquad
		V^\mu(p',p)= i e\,(p'+p)^\mu
		\]
		\item $\propto e^2$, $\mathcal L_{\text{int}}^{(e^2)}
		= 2e^2\cdot\frac12 A_\mu A^\mu\,\phi^\ast\phi$ (seagull)
		\[
		\feynmandiagram [baseline=(v.base), horizontal=a to b] {
			a [particle=\(\phi(p)\)] -- [scalar] v -- [scalar] b [particle=\(\phi(p')\)],
			v -- [photon] c [particle=\(A_\mu\)],
			v -- [photon] d [particle=\(A_\nu\)],
		};
		\qquad\Rightarrow\qquad
		V^{\mu\nu}= 2 i e^2\,g^{\mu\nu}
		\quad\text{(seagull)}
		\]
	\end{itemize}
	\end{solution}
	
	\begin{tcolorbox}[problem]
		\textbf{7.} Starting from \(\mathcal L_{\text{gf}}=-\dfrac{1}{2\xi}(\partial\cdot A)^2\), derive the photon propagator in covariant \(\xi\) gauge.
	\end{tcolorbox}
	\begin{solution}[11.7]
		We consider the following Lagrangian of QED:
		\[
		\mathcal L
		=-\frac14 F_{\mu\nu}F^{\mu\nu}-\frac{1}{2\xi}(\partial\!\cdot\!A)^2
		,\qquad
		F_{\mu\nu}=\partial_\mu A_\nu-\partial_\nu A_\mu
		\]
		Up to a total derivative, this Lagrangian can be rewritten as:
		\[
		\mathcal L
		\sim\frac12\,A_\mu
		\Bigl[g^{\mu\nu}\partial^2-\Bigl(1-\frac{1}{\xi}\Bigr)\partial^\mu\partial^\nu\Bigr]
		A_\nu
		\]
		Next, as we did in problem 3, we need to solve the following equation:
		\[
		\Bigl[g^{\mu\nu}\partial^2-\Bigl(1-\frac{1}{\xi}\Bigr)\partial^\mu\partial^\nu\Bigr]D_{\nu\rho}(x-y)=i\,\delta^\mu_{\ \rho}\,\delta^{(4)}(x-y)
		\]
		By the Fourier transformation, $\partial_\mu\to i p_\mu$. So we have:
		\[
		\left[-p^2 g^{\mu\nu}
		+\Bigl(1-\frac{1}{\xi}\Bigr)p^\mu p^\nu\right]D_{\nu\rho}(p)=i\delta^{\mu}_{\nu}
		\]
		Inverse it, one can check it like following:
\[
\begin{aligned}
	&\left[-p^2 g^{\mu\nu}
	+\Bigl(1-\frac{1}{\xi}\Bigr)p^\mu p^\nu\right]D_{\nu\rho}(p)
	=\left[-p^2 g^{\mu\nu}
	+\Bigl(1-\frac{1}{\xi}\Bigr)p^\mu p^\nu\right]
	\frac{-i}{p^2}\left(g_{\nu\rho}-(1-\xi)\frac{p_\nu p_\rho}{p^2}\right)
	\\
	=&\frac{-i}{p^2}\Biggl[
	-p^2 g^{\mu\nu}g_{\nu\rho}
	+p^2(1-\xi)g^{\mu\nu}\frac{p_\nu p_\rho}{p^2}
	+\Bigl(1-\frac{1}{\xi}\Bigr)p^\mu p^\nu g_{\nu\rho}
	-\Bigl(1-\frac{1}{\xi}\Bigr)(1-\xi)p^\mu p^\nu\frac{p_\nu p_\rho}{p^2}
	\Biggr]
	\\
	=&\frac{-i}{p^2}\Biggl[
	-p^2\delta^\mu_{\ \rho}
	+(1-\xi)p^\mu p_\rho
	+\Bigl(1-\frac{1}{\xi}\Bigr)p^\mu p_\rho
	-\Bigl(1-\frac{1}{\xi}\Bigr)(1-\xi)p^\mu p_\rho
	\Biggr]
	\\
	=&\frac{-i}{p^2}\Biggl[
	-p^2\delta^\mu_{\ \rho}
	+\Bigl((1-\xi)+\Bigl(1-\frac{1}{\xi}\Bigr)-\Bigl(1-\frac{1}{\xi}\Bigr)(1-\xi)\Bigr)p^\mu p_\rho
	\Biggr]
	\\
	=&\frac{-i}{p^2}\Biggl[
	-p^2\delta^\mu_{\ \rho}
	+\Bigl(1-\xi+1-\frac{1}{\xi}-\Bigl(1-\xi-\frac{1}{\xi}+1\Bigr)\Bigr)p^\mu p_\rho
	\Biggr]
	\\
	=&\frac{-i}{p^2}\Bigl[-p^2\delta^\mu_{\ \rho}\Bigr]
	=i\,\delta^\mu_{\ \rho}\,.
\end{aligned}
\]

		We can add the regulator $i\epsilon$ by hand. So 	\textbf{MY FINAL ANSWER IS}:
				\[
		\boxed{
			D_{\mu\nu}(p)
			=\frac{-i}{p^2+i\epsilon}
			\left(
			g_{\mu\nu}-(1-\xi)\frac{p_\mu p_\nu}{p^2}
			\right)
		}
		\]
	\end{solution}
	
	\begin{tcolorbox}[problem]
		\textbf{8.} Show explicitly that the \(\xi\)-dependent part of the photon propagator drops out of a tree-level amplitude with conserved external current \(q_\mu J^\mu=0\).
	\end{tcolorbox}
	\begin{solution}[11.8]
		The Noether current of global $U(1)$ symmetry is:
		\[
		j^\mu=i\left(\phi^*\partial^\mu\phi-(\partial^\mu\phi^*)\phi\right)
		\]
		The classical (tree-level) conservation law is:
		\[
		\partial_\mu j^\mu=0
		\]
		In the momentum space, it can be showed diagrammatically:
		\[
		\feynmandiagram [baseline=(v.base), horizontal=a to b] {
			a [particle=\(\phi(p)\)] -- [charged scalar,momentum=\(p\)] v -- [charged scalar,momentum=\(p'\)] b [particle=\(\phi(p')\)],
			v -- [photon, momentum=$q$] c [particle=\(A_\mu(q)\)],
		};
		\quad \partial_\mu J^\mu=0\quad\Rightarrow\quad q^\mu V_\mu(p,p')\sim q^\mu(p+p^{\prime})_\mu\overset{\text{OS}}{=}p^{\prime2}-p^2
		\]
		At tree level, the photon propagator always comes together with the \(A\phi\phi^*\) vertex above, and the two scalar fields are external on-shell legs. The Ward identity written above shows that, when combined with the vertex, the terms in \(D_{\mu\nu}(q)\) proportional to \(q^\mu\) do not contribute to the amplitude. We also see that the \(\xi\)-dependent term in the photon propagator is precisely \((1-\xi)\frac{q_\mu q_\nu}{q^2}\propto q^\mu\). Therefore this part (at least at tree level) gives no contribution to the amplitude, and the tree-level scattering amplitude is \(\xi\)-independent.
	\end{solution}
	
	\begin{tcolorbox}[problem]
		\textbf{9.} Derive the scalar-QED 3-point vertex rule \(ie(p+p')_\mu\) with a clear momentum-flow convention.
	\end{tcolorbox}
	\begin{solution}[11.9]
		Almost all the works have done in problem 6. There, we derived the related Lagrangian is given by:
		$$\mathcal L_{\text{int}}^{(e)}
		= i e\,A^\mu\bigl(\phi\,\partial_\mu\phi^\ast-\phi^\ast\,\partial_\mu\phi\bigr)$$
		Next, we only need to note that if $|k\rangle$ is an incoming selectron state, then $\langle0|\varphi(x)|k\rangle=e^{ikx}$ and $\langle0|\varphi^\dagger(x)|k\rangle=0;$ and if $\langle k^\prime|$ is an outgoing selectron state, then $\langle k^\prime|\varphi^\dagger(x)|0\rangle=e^{-ik^{\prime}x}$ and $\langle0|\varphi(x)|k\rangle=0.$ Therefore, in free-field theory
		\[\langle k^{\prime}|(\partial_\mu\varphi^\dagger)\varphi|k\rangle=-ik_\mu^{\prime}e^{-i(k^{\prime}-k)x},\quad\langle k^{\prime}|\varphi^\dagger\partial_\mu\varphi|k\rangle=+ik_\mu e^{-i(k^{\prime}-k)x}
		\]
		This implies that the vertex factor  is given by $i(ie)[(-ik_\mu^{\prime})-(ik_\mu)]=ie(k+k^{\prime})_\mu$. 
		
		So, \textbf{MY FINAL ANSWER IS}:
				\[
				\boxed{
		\feynmandiagram [baseline=(v.base), horizontal=a to b] {
			a [particle=\(\phi(k)\)] -- [charged scalar,momentum=\(k\)] v -- [charged scalar,momentum=\(k'\)] b [particle=\(\phi(k')\)],
			v -- [photon] c [particle=\(A_\mu(q)\)],
		};
		\qquad\Rightarrow\qquad
		V^\mu(k',k)= i e\,(k'+k)^\mu}
		\]
	\end{solution}
	
	\begin{tcolorbox}[problem]
		\textbf{10.} Consider scalar QED scattering \(\phi\phi^*\to\phi\phi^*\) via one-photon exchange.
		Compute the tree-level amplitude in general \(\xi\) gauge and show the physical result is \(\xi\)-independent (use current conservation).
	\end{tcolorbox}
	\begin{solution}[11.10]
		We have the following two diagrams:
			\[
				\feynmandiagram [baseline=(current bounding box.center),horizontal=a to b] {
					i1 [particle=$\phi$] -- [charged scalar, momentum'=$p_1$] a
					-- [charged scalar, rmomentum'=$p_2$] i2 [particle=$\phi^*$],
					a -- [photon, edge label'=$\gamma$] b,
					f1 [particle=$\phi^*$] -- [charged scalar, rmomentum'=$p_4$] b
					-- [charged scalar,momentum'=$p_3$] f2 [particle=$\phi$],
				};
				\quad 
				\begin{aligned}
					i\mathcal{M}_{s}=& ie(p_1-p_2)^\mu\times\left[\frac{-i}{(p_1+p_2)^2+i\epsilon}\left(g_{\mu\nu}-(1-\xi)\frac{(p_1+p_2)_\mu(p_1+p_2)^\nu)}{(p_1+p_2)^2}\right)\right]\\
					&\times ie(p_3-p_4)^\nu
				\end{aligned}
			\]
			\[
			% t-channel
				\feynmandiagram [baseline=(current bounding box.center),horizontal=i1 to f1, vertical=a to b] {
					i1 [particle=$\phi$] -- [charged scalar,momentum'=$p_1$] a -- [charged scalar,momentum'=$p_3$] f1 [particle=$\phi$],
					i2 [particle=$\phi^*$] -- [anti charged scalar, momentum=$p_2$] b -- [anti charged scalar,momentum=$p_4$] f2 [particle=$\phi^*$],
					a -- [photon, edge label=$\gamma$] b,
				};
				\quad
				\begin{aligned}
					i\mathcal{M}_{t}=& ie(p_1+p_3)^\mu\times\left[\frac{-i}{(p_1-p_3)^2+i\epsilon}\left(g_{\mu\nu}-(1-\xi)\frac{(p_1-p_3)_\mu(p_1-p_3)^\nu)}{(p_1-p_3)^2}\right)\right]\\
					&\times ie(p_2+p_4)^\nu
				\end{aligned}
			\]
			As we saw in the last problem, by current conservation:
			\[
			\begin{aligned}
				i\mathcal{M}_{s}&\supset (\cdots)\times (1-\xi)(p_1-p_2)\cdot(p_1+p_2) = 0\\
				i\mathcal{M}_{t}&\supset (\cdots)\times (1-\xi)(p_1+p_3)\cdot(p_1-p_3) = 0
			\end{aligned}
			\]
			Therefore we can drop out the $\xi$--dependent terms, the final amplitude indenpend on $\xi$. So, \textbf{MY FINAL ANSWER IS}:
			\[
			\begin{aligned}
				i\mathcal{M}&=i\mathcal{M}_s+\mathcal{M}_{t}\\
				&=ie^2(p_1-p_2)\cdot(p_3-p_4)\frac{1}{(p_1+p_2)^2+i\epsilon}+ie^2(p_1+p_3)\cdot(p_2+p_4)\frac{1}{(p_1-p_3)^2+i\epsilon}\\
				&=\boxed{ie^2\left(\frac{u-t}{s}+\frac{s-u}{t}\right)}
			\end{aligned}
			\]
			The last equality follows from the Mandelstam variables defined as follows.
			\[
			\begin{gathered}
				s=(p_1+p_2)^2=(p_3+p_4)^2=2p_1\cdot p_2-2m^2=2p_3\cdot p_4-2m^2\\
				t=(p_1-p_3)^2=(p_2-p_4)^2=-2p_1\cdot p_3-2m^2=-2p_2\cdot p_4-2m^2\\
				u=(p_1-p_4)^2=(p_2-p_3)^2=-2p_1\cdot p_4-2m^2=-2p_2\cdot p_3-2m^2
			\end{gathered}
			\]
			It is clear $\mathcal{M}$ is $\xi$--independent.
	\end{solution}
	
\end{document}

	\clearpage
	\documentclass[11pt]{article}

\usepackage[margin=1.6cm]{geometry}
\usepackage{amsmath,amssymb,amsthm}
\usepackage[hidelinks]{hyperref}
\usepackage[most]{tcolorbox}
\usepackage[]{cancel}
\usepackage{tikz-feynman}
\usepackage{slashed} % 用于 \not!
\usepackage{simpler-wick}   % 绘制Wick收缩
\usepackage{braket}
% --- "Solution" environment (proof-style) ---
\newenvironment{solution}[1][]{\begin{proof}[\textbf{Solution}]}{\end{proof}}
% --- Rounded box style for each problem (whole statement) ---
\tcbset{
    problem/.style={
        enhanced,
        boxrule=0.7pt,
        arc=3mm,
        left=3mm,right=3mm,top=2mm,bottom=2mm,
        colback=white,
        colframe=black,
        before skip=10pt,
        after skip=6pt,
    }
}

\begin{document}
\begin{center}
    {\LARGE \textbf{Solution Manual of Problem Set 12} \par}
    \vspace{1em}
    {\large Bufan Zheng\\ \href{mailto:bufan@g.ecc.u-tokyo.ac.jp}{\texttt{bufan@g.ecc.u-tokyo.ac.jp}} \par}
\end{center}




\begin{tcolorbox}[problem]
\textbf{1.} Verify directly from Grassmann rules that $$\int d\theta\,(a+b\theta)=b.$$
\end{tcolorbox}
\begin{solution}[12.1]
	\textbf{MY FINAL ANSWER IS}:
	\[
	\int d\theta (a+ b\theta)= a\int d\theta 1+b\int{d\theta}\theta =0+b\cdot 1=b
	\]
\end{solution}

\begin{tcolorbox}[problem]
\textbf{2.} For two Grassmann variables $\theta_1,\theta_2$, compute $$\int d\theta_2\,d\theta_1\,\theta_1\theta_2 \quad\text{and}\quad \int d\theta_2\,d\theta_1\,\theta_2\theta_1.$$ Explain the sign.
\end{tcolorbox}
\begin{solution}[12.2]
	\textbf{MY FINAL ANSWER IS}:
	\[
	\int{d\theta_2}d\theta_1\theta_1\theta_2=-\int d\theta_2 d\theta_1\theta_2\theta_1=\int d\theta_2\theta_2\int d\theta_1\theta_1 = 1\cdot 1=\boxed{1}
	\]
	
	\[
	\int{d\theta_2}d\theta_1\theta_2\theta_1=-\int d\theta_2\theta_2 d\theta_1\theta_1=-\int d\theta_2\theta_2\int d\theta_1\theta_1 = -\left(1\cdot 1\right)=\boxed{-1}
	\]
	
	They have different sign since $\theta_1\theta_2=-\theta_2\theta_1$.
\end{solution}

\begin{tcolorbox}[problem]
\textbf{3.} Prove $$\int d\psi^\dagger\,d\psi\,e^{\psi^\dagger A\psi}=\det A$$ in the $1\times 1$ and $2\times 2$ cases.
\end{tcolorbox}
\begin{solution}[12.3]
	For the $1\times 1$ case:
	\[
	\begin{aligned}
		\int d\psi^\dagger d\psi e^{\psi^\dagger A_{11}\psi}=	\int d\psi^\dagger d\psi \left(1+\psi^\dagger A_{11}\psi+0\right)
		=\int d\psi^\dagger d\psi 1+\int d\psi^\dagger d\psi \psi^\dagger A_{11}\psi=-A_{11}=\boxed{-\det A}
		\end{aligned}
	\]
		For the $2\times 2$ case:
	$$
	\begin{aligned}
		&\int d\psi^\dagger d\psi e^{\psi^\dagger A_{11}\psi}=\int d\psi_2^\dagger\,d\psi_2\,d\psi_1^\dagger\,d\psi_1\,
		\exp\!\Big(\psi_1^\dagger(A_{11}\psi_1+A_{12}\psi_2)+\psi_2^\dagger(A_{21}\psi_1+A_{22}\psi_2)\Big)\\
		&=\int d\psi_2^\dagger\,d\psi_2\,d\psi_1^\dagger\,d\psi_1\,
		\Bigg[1+\Big(\psi_1^\dagger(A_{11}\psi_1+A_{12}\psi_2)+\psi_2^\dagger(A_{21}\psi_1+A_{22}\psi_2)\Big)\\
		&\quad+\frac12\Big(\psi_1^\dagger(A_{11}\psi_1+A_{12}\psi_2)+\psi_2^\dagger(A_{21}\psi_1+A_{22}\psi_2)\Big)^2\Bigg]\\
		&=\frac12\int d\psi_2^\dagger\,d\psi_2\,d\psi_1^\dagger\,d\psi_1\,
		\Big(A_{11}\psi_1^\dagger\psi_1+A_{12}\psi_1^\dagger\psi_2+A_{21}\psi_2^\dagger\psi_1+A_{22}\psi_2^\dagger\psi_2\Big)^2\\
		&=\frac12\int d\psi_2^\dagger\,d\psi_2\,d\psi_1^\dagger\,d\psi_1\,
		\Big(2A_{11}A_{22}\,\psi_1^\dagger\psi_1\psi_2^\dagger\psi_2+2A_{12}A_{21}\,\psi_1^\dagger\psi_2\psi_2^\dagger\psi_1\Big)\\
		&=\int d\psi_2^\dagger\,d\psi_2\,d\psi_1^\dagger\,d\psi_1\,
		\Big(A_{11}A_{22}\,\psi_1^\dagger\psi_1\psi_2^\dagger\psi_2-A_{12}A_{21}\,\psi_1^\dagger\psi_1\psi_2^\dagger\psi_2\Big)\\
		&=(A_{11}A_{22}-A_{12}A_{21})\int d\psi_2^\dagger\,d\psi_2\,d\psi_1^\dagger\,d\psi_1\,\psi_1^\dagger\psi_1\psi_2^\dagger\psi_2\\
		&=(A_{11}A_{22}-A_{12}A_{21})\Big(\int d\psi_1^\dagger\,d\psi_1\,\psi_1^\dagger\psi_1\Big)\Big(\int d\psi_2^\dagger\,d\psi_2\,\psi_2^\dagger\psi_2\Big)\\
		&=(A_{11}A_{22}-A_{12}A_{21})\cdot -1\cdot -1
		=A_{11}A_{22}-A_{12}A_{21}
		=\boxed{\det A}.
	\end{aligned}
	$$
\end{solution}

\begin{tcolorbox}[problem]
\textbf{4.} For a real antisymmetric $2n\times 2n$ matrix $A$, show in the $2\times 2$ case that $$\int d^2\lambda\,\exp\!\left(\frac12\,\lambda^T A\lambda\right)=\mathrm{Pf}(A).$$
\end{tcolorbox}
\begin{solution}[12.4]
	\textbf{MY FINIAL ANSWER IS}
	\[
	\begin{aligned}
		\int d^2\lambda \exp\left(\frac{1}{2}\lambda^T A \lambda\right)
		&= \int d\theta_2 d\theta_1 \exp\left[ \frac{1}{2} \begin{pmatrix} \theta_1 & \theta_2 \end{pmatrix} \begin{pmatrix} 0 & a \\ -a & 0 \end{pmatrix} \begin{pmatrix} \theta_1 \\ \theta_2 \end{pmatrix} \right] \\
		&= \int d\theta_2 d\theta_1 \exp\left[ \frac{1}{2} \left( a\theta_1\theta_2 - a\theta_2\theta_1 \right) \right] \\
		&= \int d\theta_2 d\theta_1 \exp\left[ \frac{1}{2} \left( -a\theta_2\theta_1 - a\theta_2\theta_1 \right) \right] \quad (\because \theta_1\theta_2 = -\theta_2\theta_1) \\
		&= \int d\theta_2 d\theta_1 \exp(-a\theta_2\theta_1) \\
		&= \int d\theta_2 d\theta_1 \left( 1 - a\theta_2\theta_1 \right) \quad \\
		&= \cancelto{0}{\int d\theta_2 d\theta_1 \cdot 1} - a \int d\theta_2 d\theta_1 \theta_2 \theta_1 \\
		&= a \int d\theta_1 \left( \int d\theta_2 \theta_2 \right) \theta_1 \\
		&= a \int d\theta_1 (1) \theta_1 \\
		&= a \int d\theta_1 \theta_1 \\
		&= \mathrm{Pf}(A)
	\end{aligned}
	\]
\end{solution}

\begin{tcolorbox}[problem]
\textbf{5.} Starting from the Dirac quadratic form $\bar\psi(i\gamma\cdot\partial-m)\psi$, derive $Z[\eta,\bar\eta]=Z[0]\exp(i\bar\eta S_F\eta)$.
\end{tcolorbox}
\begin{solution}[12.5]
	\textbf{MY FINAL ANSWER IS}:
	\[
	\begin{aligned}
		Z[\eta, \bar{\eta}] 
		&= \int \mathcal{D}\bar{\psi}\mathcal{D}\psi \exp\left\{ i \int d^4x \left[ \bar{\psi}(x)(i\gamma\cdot\partial - m)\psi(x) + \bar{\eta}(x)\psi(x) + \bar{\psi}(x)\eta(x) \right] \right\} \\
		&= \int \mathcal{D}\bar{\psi}\mathcal{D}\psi \exp\left\{ i \int d^4x \left[ \left(\bar{\psi} + \bar{\eta}(i\gamma\cdot\partial - m)^{-1}\right)(i\gamma\cdot\partial - m)\left(\psi + (i\gamma\cdot\partial - m)^{-1}\eta\right) - \bar{\eta}(i\gamma\cdot\partial - m)^{-1}\eta \right] \right\} \\
		&= \underbrace{\int \mathcal{D}\bar{\psi}'\mathcal{D}\psi' \exp\left\{ i \int d^4x \, \bar{\psi}'(i\gamma\cdot\partial - m)\psi' \right\}}_{Z[0]} \cdot \exp\left\{ -i \int d^4x \, \bar{\eta}(x) (i\gamma\cdot\partial - m)^{-1} \eta(x) \right\} \\
		&= Z[0] \exp\left\{ -i \int d^4x d^4y \, \bar{\eta}(x) \left[ - S_F(x-y) \right] \eta(y) \right\} \\
		&= \boxed{Z[0] \exp\left\{ i \int d^4x di^4y \, \bar{\eta}(x) S_F(x-y) \eta(y) \right\}}
	\end{aligned}
	\]
\end{solution}
\begin{tcolorbox}[problem]
\textbf{6.} Show that $$S_F(p)=\frac{i}{\gamma\!\cdot\!p-m+i\epsilon}=\frac{i(\gamma\!\cdot\!p+m)}{p^2+m^2-i\epsilon}.$$
\end{tcolorbox}
\begin{solution}[12.6]
We need to solve the following differential equation:
\[
(i\gamma\cdot\partial-m)S_F(x-y)=i\delta^4(x-y).
\]
Using the Fourier transfornamtion, we obtain
\[
(\gamma\cdot p - m)S_F(p)=i\Rightarrow S_F(p)=\frac{i}{\gamma\cdot p-m}
\]
Next, add the $i\epsilon$ by hand, so \textbf{MY FINAL ANSWER IS}:
\[
S_{F}(p)=\frac{i}{\gamma\cdot p-m+i\epsilon}=\frac{i(\gamma\cdotp+m)}{(\gamma\cdot p-m+i\epsilon)(\gamma\cdot p+m)}=\frac{i(\gamma\cdot p+m)}{(\gamma\cdot p)^2-m^2+i\epsilon(\ldots)}=\boxed{\frac{i(\gamma\cdot p+m)}{p^2-m^2+i\epsilon}}
\]
\end{solution}
\begin{tcolorbox}[problem]
\textbf{7.} Compute $\langle 0|T\psi_\alpha(x)\bar\psi_\beta(y)|0\rangle$ from the generating functional by functional differentiation with Grassmann sources, carefully tracking signs.
\end{tcolorbox}
\begin{solution}[12.7]
we need the following formulae:
\[
\langle0|\mathrm{T}\psi_{\alpha_1}(x_1)\cdots\overline{\psi}_{\beta_1}(y_1)\cdots|0\rangle=\frac{1}{Z[0]}\frac{1}{i}\left.\frac{\delta}{\delta\overline{\eta}_{\alpha_1}(x_1)}\right.\cdots i\left.\frac{\delta}{\delta\eta_{\beta_1}(y_1)}\right.\cdots\left.Z(\overline{\eta},\eta)\right|_{\eta=\overline{\eta}=0}
\]

For the 2-pt correlator, we have:
\[
\begin{aligned}
\langle0|T\psi_\alpha(x)\bar{\psi}_\beta(y)|0\rangle&=\frac{1}{Z[0]}\frac 1i\frac{\delta}{\delta\overline{\eta}_{\alpha}(x)}\cdot i\frac{\delta}{\delta{\eta}_{\beta}(y)} Z[0] \exp\left[i\int d^4xd^4y\,\overline{\eta}(x)S_F(x-y)\eta(y)\right]\\
&=\frac{\delta}{\delta\overline{\eta}_{\alpha}(x)}\frac{\delta}{\delta{\eta}_{\beta}(y)} \left(1+i\int d^4x'd^4y'\,\bar{\eta}_{\alpha'}(x') S_F^{\alpha'\beta'}(x'-y')\eta_{\beta'}(y')\right)\\
&=- \left(i\int d^4x'd^4y'\,\frac{\delta}{\delta\overline{\eta}_{\alpha}(x)}\bar{\eta}_{\alpha'}(x') S_F^{\alpha'\beta'}(x'-y')\frac{\delta}{\delta{\eta}_{\beta}(y)}\eta_{\beta'}(y')\right)\\
&=-i{S_F}_{\alpha\beta}(x-y)
\end{aligned}
\]
\end{solution}
\begin{tcolorbox}[problem]
\textbf{8.} Using Wick's theorem for fermions, show explicitly why a closed fermion loop produces an extra minus sign.
\end{tcolorbox}
\begin{solution}[12.8]
		In complicated diagrams, one can often simplify the determination of the minus signs by noting that the product $(\bar{\psi}\psi)$, or any other pair of fermions, \textit{commutes} with any operator. Thus,
	\[
	\begin{aligned}
		\cdots 
		\wick{
			({\bar{\psi}}\c1\psi)_x(\c2{\bar{\psi}}\c3\psi)_y(\c1{\bar{\psi}}\c2\psi)_z(\c3{\bar{\psi}}\psi)_w}
		&= \cdots (+1) 
		\wick{(\bar{\psi}\c\psi)_x(\c{\bar{\psi}}}\wick{\c\psi)_z
			(\c{\bar{\psi}}}\wick{\c\psi)_y(\c{\bar{\psi}}\psi)_w}\cdots\\
		&= \cdots S_F(x - z)S_F(z - y)S_F(y - w) \cdots.
	\end{aligned}
	\]
	\noindent
	But note that in a closed loop of $n$ fermion propagators we have (e.g. $n=4$)
	\noindent
\[
\begin{aligned}
	\vcenter{\hbox{$
			\feynmandiagram [small] {
				v1 -- [fermion, quarter right] v2
				-- [fermion, quarter right] v4
				-- [fermion, quarter right] v3
				-- [fermion, quarter right] v1,
				v1 -- [scalar, dashed] e1,
				v2 -- [scalar, dashed] e2,
				v3 -- [scalar, dashed] e3,
				v4 -- [scalar, dashed] e4,
			};
			$}}\quad\quad
	&=\wick{\c2{\bar\psi}\wick{\c\psi\c{\bar\psi}}\wick{\c\psi\c{\bar\psi}}\wick{\c\psi\c{\bar\psi}}\c2{\psi}}\\
	&=(-1)\,\text{tr}\Big[
	\wick{\c\psi\c{\bar{\psi}}}\wick{\c\psi\c{\bar{\psi}}}\wick{\c\psi\c{\bar{\psi}}}\wick{\c\psi\c{\bar{\psi}}}
	\Big]\\
	&=(-1)\,\text{tr}\!\left[S_F S_F S_F S_F\right]
\end{aligned}
\]
	\noindent
	A closed fermion loop always gives a factor of $-1$ and the \textit{trace} of a product of Dirac matrices.
\end{solution}

\begin{tcolorbox}[problem]
\textbf{9.} Explain briefly but precisely how the first-order time derivative in $i\psi^\dagger\dot\psi$ is tied to the canonical anticommutator $\{\psi,\psi^\dagger\}=1$ (give the chain of logic).
\end{tcolorbox}
\begin{solution}[12.9]
\textbf{MY FINAL ANSWER IS}:
$$
\begin{aligned}
	&L(\psi,\psi^\dagger)\supset i\psi^\dagger \dot\psi\Longrightarrow\ \pi_\psi:=\frac{\partial L}{\partial\dot\psi}=i\psi^\dagger,\ \pi_{\psi^\dagger}:=\frac{\partial L}{\partial\dot\psi^\dagger}=0\\
	&\Longrightarrow\ \{\psi,\pi_\psi\}_P=1,\ \{\psi^\dagger,\pi_{\psi^\dagger}\}_P=1\ \text{(classical Poisson bracket)}\\
	&\Longrightarrow\ \text{quantize by the rule }\{\cdot,\cdot\}_P\mapsto\frac{1}{i}\{\cdot,\cdot\}_D\ \Longrightarrow\ \boxed{\{\psi,\psi^\dagger\}_D=1}.
\end{aligned}
$$
\end{solution}

\begin{tcolorbox}[problem]
\textbf{10.} In Euclidean time, derive the anti-periodic boundary condition for fermions at finite temperature and show why the Matsubara frequencies are $\omega_n=(2n+1)\pi/\beta$.
\end{tcolorbox}
\begin{solution}[12.10]
	Consider a general operator $\hat{A}$,
	$$\hat A|i\rangle=|j\rangle A_{ji},$$
	where $i,j$ run over $\downarrow$ and $\uparrow$. Then by expanding $\hat{A}|\psi\rangle$ in terms of the basis states one finds
	$$\int d\psi\:\langle\psi|\hat{A}|\psi\rangle=A_{\downarrow\downarrow}-A_{\uparrow\uparrow}=\mathrm{Tr}\Big[(-1)^{\hat{F}}\hat{A}\Big]$$
	where the fermion number $\hat{F}$ has been defined to be even for $|\downarrow\rangle$ and odd
	for $|\uparrow\rangle$. We used $\braket{\psi|\downarrow}=\psi,\braket{\psi|\uparrow}=-1$. In particular,
	$$\int d\psi\left\langle\psi,T|\psi,0\right\rangle=\mathrm{Tr}\left[(-1)^{\hat{F}}\exp(-i\hat{H}T)\right]\:.$$
	In contrast to the bosonic results, the \textit{periodic} pathintegral gives the trace weighted by $(-1)^{\hat{F}}$ and the \textit{antiperiodic} path integral gives the simple trace. Here, what we need is $\operatorname{Tr}(e^{-\beta \hat H})$, i.e. the simple trace, so we conclude that the boundary condition for fermions should be anti-periodic to obtain the simple trace.
	
	To obtain the Matsubara frequencies, we need to Fourier transform the $\psi$:
	$$
			\psi(\tau,x)=\frac1{\sqrt{\beta}}\sum_{n\in\mathbb Z}\psi_n(x)e^{-i\omega_n\tau}
	$$
	Next, by the boundary condition, \textbf{MY FINAL ANSWER IS}:
	$$
	\begin{aligned}
	&\psi(\tau+\beta,x)=-\psi(\tau,x)
		\Longrightarrow\ e^{-i\omega_n(\tau+\beta)}=-e^{-i\omega_n\tau}
		\ \Longrightarrow\ e^{-i\omega_n\beta}=-1\\&
		\ \Longrightarrow\ \omega_n\beta=(2n+1)\pi
		\ \Longrightarrow\ \omega_n=\boxed{\frac{(2n+1)\pi}{\beta}}
	\end{aligned}
	$$
	Here, $n\in\mathbb Z$.
\end{solution}

\end{document}

	\clearpage
	\documentclass[11pt]{article}

\usepackage[margin=1.6cm]{geometry}
\usepackage{amsmath,amssymb,amsthm}
\usepackage[hidelinks]{hyperref}
\usepackage[most]{tcolorbox}
\usepackage[]{cancel}
\usepackage{tikz-feynman}
\tikzfeynmanset{compat=1.1.0}
% --- "Solution" environment (proof-style) ---
\newenvironment{solution}[1][]{\begin{proof}[\textbf{Solution}]}{\end{proof}}
% --- Rounded box style for each problem (whole statement) ---
\tcbset{
    problem/.style={
        enhanced,
        boxrule=0.7pt,
        arc=3mm,
        left=3mm,right=3mm,top=2mm,bottom=2mm,
        colback=white,
        colframe=black,
        before skip=10pt,
        after skip=6pt,
    }
}

\begin{document}
\begin{center}
    {\LARGE \textbf{Solution Manual of Problem Set 13} \par}
    \vspace{1em}
    {\large Bufan Zheng\\ \href{mailto:bufan@g.ecc.u-tokyo.ac.jp}{\texttt{bufan@g.ecc.u-tokyo.ac.jp}} \par}
\end{center}


\begin{tcolorbox}[problem]
\textbf{1.} Starting from $\{\gamma^\mu,\gamma^\nu\}=2g^{\mu\nu}$, show that $[\gamma^\mu,\gamma^\nu]$ generates a representation of the Lorentz algebra on spinors via $\Sigma^{\mu\nu}=\frac14[\gamma^\mu,\gamma^\nu]$.
\end{tcolorbox}
\begin{solution}[13.1]
	\textbf{MY FINAL ANSWER IS}:
$$
\begin{aligned}
	[\Sigma^{\mu\nu},\Sigma^{\rho\sigma}]
	&=\Big[\tfrac14[\gamma^\mu,\gamma^\nu],\tfrac14[\gamma^\rho,\gamma^\sigma]\Big]
	=\tfrac1{16}[[\gamma^\mu,\gamma^\nu],[\gamma^\rho,\gamma^\sigma]]\\
	&=\tfrac1{16}\Big([[\,\gamma^\mu,\gamma^\nu\,],\gamma^\rho]\gamma^\sigma+\gamma^\rho[[\,\gamma^\mu,\gamma^\nu\,],\gamma^\sigma]-[[\,\gamma^\mu,\gamma^\nu\,],\gamma^\sigma]\gamma^\rho-\gamma^\sigma[[\,\gamma^\mu,\gamma^\nu\,],\gamma^\rho]\Big)\\
	&=\tfrac1{16}\Big([\;[[\,\gamma^\mu,\gamma^\nu\,],\gamma^\rho]\;,\gamma^\sigma]-[\;[[\,\gamma^\mu,\gamma^\nu\,],\gamma^\sigma]\;,\gamma^\rho]\Big)\\
	&=\tfrac1{16}\Big([\;\gamma^\mu\gamma^\nu\gamma^\rho-\gamma^\nu\gamma^\mu\gamma^\rho-\gamma^\rho\gamma^\mu\gamma^\nu+\gamma^\rho\gamma^\nu\gamma^\mu\;,\gamma^\sigma]-[\;\gamma^\mu\gamma^\nu\gamma^\sigma-\gamma^\nu\gamma^\mu\gamma^\sigma-\gamma^\sigma\gamma^\mu\gamma^\nu+\gamma^\sigma\gamma^\nu\gamma^\mu\;,\gamma^\rho]\Big)\\
	&=\tfrac1{16}\Big([\;\gamma^\mu(2g^{\nu\rho}-\gamma^\rho\gamma^\nu)-\gamma^\nu(2g^{\mu\rho}-\gamma^\rho\gamma^\mu)-(2g^{\rho\mu}-\gamma^\mu\gamma^\rho)\gamma^\nu+(2g^{\rho\nu}-\gamma^\nu\gamma^\rho)\gamma^\mu\;,\gamma^\sigma]\\
	&\qquad\qquad-[\;\gamma^\mu(2g^{\nu\sigma}-\gamma^\sigma\gamma^\nu)-\gamma^\nu(2g^{\mu\sigma}-\gamma^\sigma\gamma^\mu)-(2g^{\sigma\mu}-\gamma^\mu\gamma^\sigma)\gamma^\nu+(2g^{\sigma\nu}-\gamma^\nu\gamma^\sigma)\gamma^\mu\;,\gamma^\rho]\Big)\\
	&=\tfrac1{16}\Big([\;4(g^{\nu\rho}\gamma^\mu-g^{\mu\rho}\gamma^\nu)\;,\gamma^\sigma]-[\;4(g^{\nu\sigma}\gamma^\mu-g^{\mu\sigma}\gamma^\nu)\;,\gamma^\rho]\Big)\\
	&=\tfrac14\Big(g^{\nu\rho}[\gamma^\mu,\gamma^\sigma]-g^{\mu\rho}[\gamma^\nu,\gamma^\sigma]-g^{\nu\sigma}[\gamma^\mu,\gamma^\rho]+g^{\mu\sigma}[\gamma^\nu,\gamma^\rho]\Big)\\
	&=g^{\nu\rho}\Sigma^{\mu\sigma}-g^{\mu\rho}\Sigma^{\nu\sigma}-g^{\nu\sigma}\Sigma^{\mu\rho}+g^{\mu\sigma}\Sigma^{\nu\rho}.
\end{aligned}
$$
It is clear that this is the commutator in Lorentz algebra.
\end{solution}

\begin{tcolorbox}[problem]
\textbf{2.} Show that the Dirac Lagrangian $\bar\psi(i\gamma_\mu\partial^\mu-m)\psi$ is real (Hermitian action) given appropriate hermiticity properties of the $\gamma^\mu$.
\end{tcolorbox}
\begin{solution}[13.2]
	$$
	\begin{aligned}\left(\bar{\psi}\left(i\gamma_\mu\partial^\mu-m\right)\psi\right)^\dagger&=\left(i\left.\psi^\dagger\gamma^0\gamma_\mu\partial^\mu\psi-m\right.\psi^\dagger\gamma^0\psi\right)^{\dagger}\\&=-i\left(\partial^\mu\psi\right)^\dagger(\gamma_\mu)^\dagger(\gamma^0)^\dagger\psi-m\psi^\dagger(\gamma^0)\\&=-i\left.\partial^{\mu}\psi^{\dagger}(\gamma_{\mu})^{\dagger}\gamma^{0}\psi\right.-\left.m\psi^{\dagger}\gamma^{0}\psi\right.\\&=-i\partial^\mu\left(\psi^\dagger(\gamma_\mu)^\dagger\gamma^0\psi\right)+i\psi^\dagger(\gamma_\mu)^\dagger\gamma^0\partial^\mu\psi-m\psi^\dagger\gamma^0\psi
\end{aligned}
$$
		If we have $(\gamma_\mu)^\dagger=\gamma^0\gamma_\mu\gamma^0$, then
	$$
		\begin{aligned}
		\text{LHS}&=-\partial^{\mu}{\left(i\left.\psi^{\dagger}(\gamma_{\mu})^{\dagger}\gamma^{0}\psi\right)\right.}+\left.\psi^{\dagger}\gamma^{0}(i\gamma_{\mu}\partial^{\mu}-m)\psi\right.\\&=-\partial^\mu\left(i\bar{\psi}\gamma_\mu\psi\right)+\bar{\psi}(i\gamma_\mu\partial^\mu-m)\psi\\&=\bar{\psi}\left(i\gamma_\mu\partial^\mu-m\right)\psi-\partial^\mu\left(i\bar{\psi}\gamma_\mu\psi\right).\end{aligned}
	$$
	i.e. the Lagrangian is real up to a total derivative term. Therefore, we need the following hermiticity properties of the $\gamma^\mu$:
	$$
	\boxed{
	(\gamma_\mu)^\dagger=\gamma^0\gamma_\mu\gamma^0\iff\bar{\gamma}_\mu=\gamma_\mu	}
	$$
\end{solution}
\begin{tcolorbox}[problem]
\textbf{3.} In even $D$, define $\gamma_\ast$ and show that $P_{L/R}=\frac12(1\mp\gamma_\ast)$ are projectors.
\end{tcolorbox}
\begin{solution}[13.3] Fist, we can choose the following ansatz:
	\[
	\gamma_*\propto\gamma^0\gamma^1\cdots\gamma^{D-1}
	\]
	Next, we check some properties of $\gamma_*$:
	$$
	\begin{aligned}
		\{\gamma_*,\gamma^\mu\}
		&=\gamma_*\gamma^\mu+\gamma^\mu\gamma_*
		=(\gamma^0\gamma^1\cdots\gamma^{D-1})\gamma^\mu+\gamma^\mu(\gamma^0\gamma^1\cdots\gamma^{D-1})\\
		&=(\gamma^0\cdots\gamma^{D-1})\gamma^\mu+(-1)^{D-1}\gamma^0\cdots\gamma^{D-1}\gamma^\mu
		=\bigl(1+(-1)^{D-1}\bigr)\gamma_*\gamma^\mu
		=0\qquad(D\ \text{even}).
	\end{aligned}
	$$
	$$
		\gamma_*^2
		=(\gamma^0\gamma^1\cdots\gamma^{D-1})(\gamma^0\gamma^1\cdots\gamma^{D-1})
		=(-1)^{\frac{D(D-1)}{2}}\prod_{\mu=0}^{D-1}(\gamma^\mu)^2\\
		=(-1)^{\frac{D(D-1)}{2}}\prod_{\mu=0}^{D-1}g^{\mu\mu}\,\mathbf 1=(-1)^{\frac{D(D-1)}{2}+1}\mathbf{1}
	$$
	we can choose the normalization such that $\gamma_*^2=1$. So, \textbf{MY FINAL ANSWER IS}
	$$
	\boxed{
	\gamma_*=(-1)^{\frac{D(D-1)}{2}+1}\gamma^0\gamma^1\cdots\gamma^{D-1}
}
	$$
	Next, we check $P_{L/R}^2=P_{L/R}$
	$$
	\begin{aligned}
		P_L^2
		&=\Big(\tfrac12(1-\gamma_*)\Big)^2
		=\tfrac14(1-2\gamma_*+\gamma_*^2)
		=\tfrac14(1-2\gamma_*+1)
		=\tfrac12(1-\gamma_*)
		=P_L.
	\end{aligned}
	$$
	
	$$
	\begin{aligned}
		P_R^2
		&=\Big(\tfrac12(1+\gamma_*)\Big)^2
		=\tfrac14(1+2\gamma_*+\gamma_*^2)
		=\tfrac14(1+2\gamma_*+1)
		=\tfrac12(1+\gamma_*)
		=P_R.
	\end{aligned}
	$$
	
\end{solution}
\begin{tcolorbox}[problem]
\textbf{4.} State clearly what it means for a spinor representation to be real, complex, or quaternionic (pseudoreal) in this context.
\end{tcolorbox}
\begin{solution}[13.4]
	Using the charge conjugation matrix \(C\), define \(\psi^c \equiv C\bar\psi^{\,T}\).
	The spinor representation can be classified by whether one can impose a reality condition relating \(\psi\) and \(\psi^c\),
	and by the sign obtained after applying charge conjugation twice.
	\begin{itemize}
		\item 	\textbf{Real (Majorana).}
		One can consistently impose the Majorana condition \(\psi^c=\psi\), with the consistency requirement
		\((\psi^c)^c=+\psi\).
		\item	\textbf{Complex.}
		The representation is not equivalent to its complex conjugate, so there is no compatible charge conjugation reality structure.
		Hence no constraint of the form \(\psi^c=\pm\psi\) can be imposed, and \(\psi\) must be a Dirac spinor.
		\item 	\textbf{Quaternionic / Pseudoreal.}
		For a single spinor one has \((\psi^c)^c=-\psi\), so \(\psi^c=\psi\) is inconsistent.
		Introducing an \(SU(2)\simeq Sp(1)\) doublet \(\psi^i\,(i=1,2)\), one can impose the symplectic Majorana condition
		\((\psi^i)^c=\epsilon^{ij}\psi_j\) with \(\epsilon^{12}=1\), which is consistent and satisfies \(((\psi^i)^c)^c=+\psi^i\). (Problem 9 will give you an example)
	\end{itemize}
\end{solution}

\begin{tcolorbox}[problem]
\textbf{5.} Using the rule $r=p-q\;(\mathrm{mod}\ 8)$, compute $r$ for Lorentzian mostly-plus signature in $D=2,3,4,6,10$.
\end{tcolorbox}
\begin{solution}[13.5]
	mostly-plus signature means $(p,q)=(D-1,1)$, therefore, we have $r= D-2\quad(\mathrm{mod~}8)$. \textbf{MY FINAL ANSWER IS}
	$$
	\boxed{
	\begin{aligned}D=2:&\quad r\equiv0\quad(\mathrm{mod~}8),\\D=3:&\quad r\equiv1\quad(\mathrm{mod~}8),\\D=4:&\quad r\equiv2\quad(\mathrm{mod~}8),\\D=6:&\quad r\equiv4\quad(\mathrm{mod~}8),\\D=10:&\quad r\equiv8\equiv0\quad(\mathrm{mod~}8).\end{aligned}
}
	$$
	
\end{solution}

\begin{tcolorbox}[problem]
\textbf{6.} For each of the above $D$, state whether Majorana and/or Weyl conditions can be imposed.
\end{tcolorbox}
\begin{solution}[13.6]
Here, we conclude the condition for Weyl and Majorana.

\begin{itemize}
	\item \textbf{Weyl (W):}A Weyl (chiral) condition exists iff $D$ is even.
	\item\textbf{Majorana (M):}A (real) Majorana condition exists iff $r\equiv 0,1,2(\mathrm{mod}\ 8)$.
	\item \textbf{Majorana--Weyl (MW):} MW requires Weyl plus a compatible reality condition, in Lorentzian signature this happens for $r\equiv 0\ (\mathrm{mod}\ 8)$.
\end{itemize}
	
	\[
	\begin{array}{c|c|c|c|c}
		D & r\equiv D-2\ (\mathrm{mod}\ 8) & \mathrm{M} & \mathrm{W} & \mathrm{MW}\\
		\hline
		2  & 0 & \text{Yes} & \text{Yes} & \text{Yes}\\
		3  & 1 & \text{Yes} & \text{No}  & \text{No}\\
		4  & 2 & \text{Yes} & \text{Yes} & \text{No}\\
		6  & 4 & \text{No} & \text{Yes} & \text{No}\\
		10 & 0 & \text{Yes} & \text{Yes} & \text{Yes}
	\end{array}
	\]
	
\end{solution}
\begin{tcolorbox}[problem]
\textbf{7.} In $D=4$, explain why charge conjugation flips chirality and why this prevents a purely chiral Majorana fermion.
\end{tcolorbox}
\begin{solution}[13.7]
	For \(D=4\) Lorentzian, \(r=D-2=2\!\!\pmod{8}\), so the spinor representation is of real type and a Majorana condition can be imposed on a four-component spinor. Since \(D\) is even, one can also impose the Weyl (chirality) condition. However, these two conditions are not simultaneously compatible for a \emph{single} chiral spinor, because charge conjugation flips chirality. Concretely, with \(\psi^c \equiv C\bar\psi^{\,T}\) and chiral projectors \(P_{L,R}=\tfrac12(1\mp\gamma_5)\), one has \((P_L\psi)^c = P_R\psi^c\), hence \((\psi_L)^c\) is right-handed. Therefore the Majorana condition \(\psi^c=\psi\) cannot be imposed within a purely left-handed (or purely right-handed) Weyl subspace, it necessarily relates left- and right-handed components. This is why there is no Majorana--Weyl spinor in \(D=4\) Lorentzian signature.
\end{solution}

\begin{tcolorbox}[problem]
\textbf{8.} In $D=10$, explain why Majorana--Weyl is possible.
\end{tcolorbox}
\begin{solution}[13.8]
	For $D=10:$
	$$r=D-2=8\equiv0\pmod8,$$
	It means in 10 dimension, the Clifford algebra is a real type one, hence Majorana possible, and $D$ even means Weyl possible. In
	this case, one can impose both:
	$$\psi=\psi^c,\quad\gamma_*\psi=\pm\psi,$$
	so Majorana-Weyl spinors exist.
\end{solution}

\begin{tcolorbox}[problem]
\textbf{9.} Give one example of a ``symplectic Majorana'' condition (conceptually), what extra index structure is needed and why?
\end{tcolorbox}
\begin{solution}[13.9]
	In \(D=6\) Lorentzian signature,
	define charge conjugation by \(\psi^c \equiv C\bar\psi^{\,T}\).
	In this case one has \((\psi^c)^c=-\psi\), so imposing the usual Majorana condition \(\psi^c=\psi\) would immediately give
	\(\psi=(\psi^c)^c=-\psi\), hence \(\psi=0\).
	Therefore a single Majorana spinor does not exist.
	
	The remedy is to introduce an additional \(SU(2)\simeq Sp(1)\) index so that the fermions come as a doublet
	\(\psi^i\,(i=1,2)\), and impose the symplectic Majorana condition
	\((\psi^i)^c=\epsilon^{ij}\psi_j\) with \(\epsilon^{12}=1\).
	The symplectic tensor \(\epsilon^{ij}\) supplies the extra minus sign needed so that applying charge conjugation twice gives back
	\(+\psi^i\), making the constraint consistent.
	
\end{solution}
\begin{tcolorbox}[problem]
\textbf{10.} Explain why Euclidean ``Majorana'' is subtle, what goes wrong with na\"ively demanding $\psi^\ast=\psi$ for Euclidean gamma matrices, and what is the usual practical workaround in path integrals?
\end{tcolorbox}
\begin{solution}
	
\end{solution}

\end{document}

	\clearpage
	\documentclass[11pt]{article}

\usepackage[margin=1.6cm]{geometry}
\usepackage{amsmath,amssymb,amsthm}
\usepackage[hidelinks]{hyperref}
\usepackage[most]{tcolorbox}
\usepackage[]{cancel}
\usepackage{tikz-feynman}
\usepackage{slashed}
\tikzfeynmanset{compat=1.1.0}
\newenvironment{solution}[1][]{\begin{proof}[\textbf{Solution}]}{\end{proof}}
% --- Rounded box style for each problem (whole statement) ---
\tcbset{
    problem/.style={
        enhanced,
        boxrule=0.7pt,
        arc=3mm,
        left=3mm,right=3mm,top=2mm,bottom=2mm,
        colback=white,
        colframe=black,
        before skip=10pt,
        after skip=6pt,
    }
}



\begin{document}
\begin{center}
    {\LARGE \textbf{Solution Manual of Problem Set 14} \par}
    \vspace{1em}
    {\large Bufan Zheng\\ \href{mailto:bufan@g.ecc.u-tokyo.ac.jp}{\texttt{bufan@g.ecc.u-tokyo.ac.jp}} \par}
\end{center}



\begin{tcolorbox}[problem]
\textbf{1.} Draw the two tree-level QED diagrams contributing to Compton scattering $e^-\gamma\to e^-\gamma$, label momenta, and write down the amplitude as a sum of two terms.
\end{tcolorbox}
\begin{solution}[14.1]
	We have the following Feynman diagrams:
	
\begin{figure}[h]
	\centering
	
	% s-channel
	\begin{tikzpicture}[baseline=(current bounding box.center)]
		\begin{feynman}
			\diagram [horizontal=a to b] {
				i1 [particle=$e^-(p)$] -- [fermion, momentum'=$p$] a
				-- [fermion, momentum=$p+k$] b
				-- [fermion, momentum=$p'$] f1 [particle=$e^-(p')$],
				i2 [particle={$\gamma(k,\epsilon)$}] -- [photon, momentum=$k$] a,
				b -- [photon, momentum'=$k'$] f2 [particle={$\gamma(k',\epsilon')$}],
			};
		\end{feynman}
	\end{tikzpicture}
	\hspace{1.2cm}
	% u-channel
	\begin{tikzpicture}[baseline=(current bounding box.center)]
		\begin{feynman}
			\diagram [vertical=a to b] {
				i1 [particle=$e^-(p)$] -- [fermion, momentum'=$p$] a,
				a -- [fermion, momentum=$p-k'$] b,
				b -- [fermion, momentum=$p'$] f1 [particle=$e^-(p')$],
				a -- [photon, momentum'=$k'$] f2 [particle={$\gamma(k',\epsilon')$}],
				i2 [particle={$\gamma(k,\epsilon)$}] -- [photon, momentum=$k$] b,
			};
		\end{feynman}
	\end{tikzpicture}
	
\end{figure}

	By Feynman rules, \textbf{MY FINAL ANSWER IS}:
		\[
		\boxed{
		\begin{gathered}
		i\mathcal M_s
		=(-ie)^2\,\bar u(p')\,
		\slashed\epsilon^{\,*}(k')\,
		\frac{i(\slashed p+\slashed k+m)}{(p+k)^2-m^2+i\varepsilon}\,
		\slashed\epsilon(k)\,
		u(p),\\
		i\mathcal M_u
		=(-ie)^2\,\bar u(p')\,
		\slashed\epsilon(k)\,
		\frac{i(\slashed p-\slashed k'+m)}{(p-k')^2-m^2+i\varepsilon}\,
		\slashed\epsilon^{\,*}(k')\,
		u(p).\\
	i\mathcal M = i\mathcal M_s + i\mathcal M_u,
			\end{gathered}}
	\]
\end{solution}



\begin{tcolorbox}[problem]
\textbf{2.} Using on-shell conditions, show that the tree amplitude is gauge invariant, replacing $\epsilon^\mu(k)\to k^\mu$ gives zero (Ward identity at tree level).
\end{tcolorbox}
\begin{solution}[14.2]
	We need to prove that $i\mathcal{M}({\epsilon\to k}=0)$. Since the external particles are on-shell, so we can use the following Dirac equations:
	
	% (1) External Dirac equations (on-shell spinors)
	\[
	(\slashed p-m)\,u(p)=0,
	\qquad
	\bar u(p')\,(\slashed p'-m)=0.
	\]
	\textbf{MY FINAL ANSWER IS}:
	% 1) 把入射光子的极化替换为动量, \slashed\epsilon(k)\to\slashed k
	\begin{align*}
		i\mathcal M\Big|_{\slashed\epsilon(k)\to\slashed k}
		&=(-ie)^2\,\bar u(p')\Bigg[
		\slashed\epsilon^{\,*}(k')\,
		\frac{i(\slashed p+\slashed k+m)}{(p+k)^2-m^2+i\varepsilon}\,
		\slashed k
		+\slashed k\,
		\frac{i(\slashed p-\slashed k'+m)}{(p-k')^2-m^2+i\varepsilon}\,
		\slashed\epsilon^{\,*}(k')
		\Bigg]u(p)
		\\
		&=(-ie)^2\,\bar u(p')\Bigg[
		\slashed\epsilon^{\,*}(k')\,
		\frac{i(\slashed p+\slashed k+m)}{(p+k)^2-m^2+i\varepsilon}\,
		\Big((\slashed p+\slashed k-m)-(\slashed p-m)\Big)
		\\
		&\qquad\qquad
		+\Big((\slashed p'-m)-(\slashed p-\slashed k'-m)\Big)\,
		\frac{i(\slashed p-\slashed k'+m)}{(p-k')^2-m^2+i\varepsilon}\,
		\slashed\epsilon^{\,*}(k')
		\Bigg]u(p)
		\\
		&=(-ie)^2\,\bar u(p')\Bigg[
		\slashed\epsilon^{\,*}(k')\,
		\frac{i(\slashed p+\slashed k+m)(\slashed p+\slashed k-m)}{(p+k)^2-m^2+i\varepsilon}
		-(\slashed p-\slashed k'-m)\,
		\frac{i(\slashed p-\slashed k'+m)}{(p-k')^2-m^2+i\varepsilon}\,
		\slashed\epsilon^{\,*}(k')
		\Bigg]u(p)
		\\
		&=(-ie)^2\,\bar u(p')\Bigg[
		\slashed\epsilon^{\,*}(k')\,
		\frac{i\big((p+k)^2-m^2\big)}{(p+k)^2-m^2+i\varepsilon}
		-\frac{i\big((p-k')^2-m^2\big)}{(p-k')^2-m^2+i\varepsilon}\,
		\slashed\epsilon^{\,*}(k')
		\Bigg]u(p)
		\\
		&=(-ie)^2\,i\,\bar u(p')\slashed\epsilon^{\,*}(k')u(p)
		-(-ie)^2\,i\,\bar u(p')\slashed\epsilon^{\,*}(k')u(p)
		\\
		&=0.
	\end{align*}
	% 2) 把出射光子的极化替换为动量, \slashed\epsilon^{\,*}(k')\to\slashed k'
	
	\begin{align*}
		i\mathcal M\Big|_{\slashed\epsilon^{\,*}(k')\to\slashed k'}
		&=(-ie)^2\,\bar u(p')\Bigg[
		\slashed k'\,
		\frac{i(\slashed p+\slashed k+m)}{(p+k)^2-m^2+i\varepsilon}\,
		\slashed\epsilon(k)
		+\slashed\epsilon(k)\,
		\frac{i(\slashed p-\slashed k'+m)}{(p-k')^2-m^2+i\varepsilon}\,
		\slashed k'
		\Bigg]u(p)
		\\
		&=(-ie)^2\,\bar u(p')\Bigg[
		\Big((\slashed p+\slashed k-m)-(\slashed p'-m)\Big)\,
		\frac{i(\slashed p+\slashed k+m)}{(p+k)^2-m^2+i\varepsilon}\,
		\slashed\epsilon(k)
		\\
		&\qquad\qquad
		+\slashed\epsilon(k)\,
		\frac{i(\slashed p-\slashed k'+m)}{(p-k')^2-m^2+i\varepsilon}\,
		\Big((\slashed p-m)-(\slashed p-\slashed k'-m)\Big)
		\Bigg]u(p)
		\\
		&=(-ie)^2\,\bar u(p')\Bigg[
		\frac{i(\slashed p+\slashed k-m)(\slashed p+\slashed k+m)}{(p+k)^2-m^2+i\varepsilon}\,
		\slashed\epsilon(k)
		-\slashed\epsilon(k)\,
		\frac{i(\slashed p-\slashed k'+m)(\slashed p-\slashed k'-m)}{(p-k')^2-m^2+i\varepsilon}
		\Bigg]u(p)
		\\
		&=(-ie)^2\,\bar u(p')\Bigg[
		\frac{i\big((p+k)^2-m^2\big)}{(p+k)^2-m^2+i\varepsilon}\,
		\slashed\epsilon(k)
		-\slashed\epsilon(k)\,
		\frac{i\big((p-k')^2-m^2\big)}{(p-k')^2-m^2+i\varepsilon}
		\Bigg]u(p)
		\\
		&=(-ie)^2\,i\,\bar u(p')\slashed\epsilon(k)u(p)
		-(-ie)^2\,i\,\bar u(p')\slashed\epsilon(k)u(p)
		\\
		&=0.
	\end{align*}
\end{solution}

\begin{tcolorbox}[problem]
\textbf{3.} In covariant $\xi$ gauge, show explicitly that the amplitude is independent of $\xi$ (at tree level) because external photons appear only through polarization vectors.
\end{tcolorbox}
\begin{solution}[14.3]
	In covariant $\xi$ gauge the gauge-fixing term is
	\[
	\mathcal L_{\text{gf}}=-\frac{1}{2\xi}\,(\partial_\mu A^\mu)^2,
	\]
	which modifies only the photon propagator,
	\[
	D_{\mu\nu}(q)=\frac{-i}{q^2+i\varepsilon}\left(\eta_{\mu\nu}-(1-\xi)\frac{q_\mu q_\nu}{q^2+i\varepsilon}\right).
	\]
	Hence any $\xi$-dependence of an amplitude can only come from internal photon lines.
	
	For tree-level Compton scattering $e^-\gamma\to e^-\gamma$, the two diagrams contain
	only the QED vertices $(-ie\gamma^\mu)$ and internal \emph{electron} propagators
	\[
	S_F(q)=\frac{i(\slashed q+m)}{q^2-m^2+i\varepsilon},
	\]
	which is $\xi$--independent. While the photons are \emph{external} and enter only through polarization vectors
	$\epsilon^\mu(k)$ and $\epsilon^{*\nu}(k')$. So we only need to care about the following sub-diagram:
	\[
	\begin{tikzpicture}[baseline=(current bounding box.center)]
		\begin{feynman}
			\diagram [horizontal=a to f] {
				i1 [particle=$e^-(p)$] -- [fermion, momentum'=$p$] a
				-- [fermion, momentum'=$p'$] i2 [particle=$e^-(p')$],
				a -- [photon, momentum'=$p-p'$] f [particle={$\gamma(q',\epsilon')$}],
			};
		\end{feynman}
	\end{tikzpicture}
	\quad\supset\quad q_\mu \bar u(p')\gamma^\mu u(p)
	\]
	
	\[
	\begin{aligned}
		q_\mu \bar u(p')\gamma^\mu u(p)
		&=\bar u(p')\,\slashed q\,u(p)\\
		&=\bar u(p')(\slashed p'-\slashed p)u(p)\\
		&=\bar u(p')\big[(\slashed p'-m)-(\slashed p-m)\big]u(p)\\
		&=\bar u(p')(\slashed p'-m)u(p)-\bar u(p')(\slashed p-m)u(p)\\
		&=0.
	\end{aligned}
	\]
	Therefore, any term in $D_{\mu\nu}$ that is proportional to the momentum $k_\mu$ does not contribute to the amplitude. Since the $\xi$--dependent part of $D_{\mu\nu}$ is precisely of this form, the final amplitude is $\xi$--independent.
\end{solution}

\begin{tcolorbox}[problem]
\textbf{4.} In $D$ dimensions, show that a physical photon has $D-2$ polarizations. Give a simple argument in momentum space.
\end{tcolorbox}
\begin{solution}[14.4]
	The polarization vector of a photon is constrained by the following on-shell equation:
	\[
	k^\mu \epsilon_\mu = 0,
	\]
	which reduces the degrees of freedom to $D-1$. Due to gauge invariance:
	\[
	\epsilon^\mu \sim \epsilon^\mu + k^\mu,
	\]
	the physical degrees of freedom of the photon are eventually reduced to $D-2$.
\end{solution}
\begin{tcolorbox}[problem]
\textbf{5.} Write the polarization sum for external photons in a way consistent with gauge invariance (i.e.\ explain why you can effectively use $-\eta^{\mu\nu}$ inside gauge-invariant expressions).
\end{tcolorbox}
\begin{solution}[14.5]
consider an arbitrary QED process involving an external photon with momentum $k$:
\[
\begin{tikzpicture}[baseline=(current bounding box.center)]
	\begin{feynman}
		% 1. 定义中心标准 blob
		% 可以通过 minimum size 稍微调整它的大小
		\vertex [blob, minimum size=1cm] (c) {};
		
		% 2. 定义周围的顶点位置（手动指定以模拟原图的随意感）
		% 左侧入射
		\vertex [above left=2cm of c] (i1);
		\vertex [left=1.8cm of c, yshift=0.5cm] (i2);
		\vertex [below left=2.2cm of c, yshift=1cm] (i3);
		\vertex [below left=1cm of c,yshift=-1.2cm, rotate=-20] (i4);
		
		% 右侧出射
		% 关键的带动量 k 的光子顶点位置
		\vertex [above right=2.5cm of c] (o_k); 
		% 其他出射粒子
		\vertex [below right=2cm of c] (o2);
		\vertex [right=1.8cm of c, yshift=-0.5cm] (o3);
		
		% 3. 绘制连接线 (使用 \diagram* 手动模式)
		\diagram* {
			% 入射粒子 (混合直线和波浪线)
			(i1) -- [fermion] (c),
			(i2) -- [photon] (c),
			(i3) -- [photon] (c),
			(i4) -- [fermion] (c),
			
			% 关键：带有动量 k 的出射光子
			% momentum 选项会自动添加箭头和标签
			(c) -- [photon, momentum=\(k\)] (o_k),
			
			% 其他出射粒子
			(c) -- [fermion] (o2),
			(c) -- [fermion] (o3),
		};
	\end{feynman}
\end{tikzpicture}
	\qquad= i\mathcal{M}(k)\equiv i\mathcal{M}^\mu(k)\epsilon_\mu^*(k)
\]
We can choose the reference frame such that $k^\mu=(k,0,0,k)$. Hence the polarization vector is:
\[
\epsilon_1^\mu=(0,1,0,0);\quad\epsilon_2^\mu=(0,0,1,0).
\]
With these conventions, we have
\[
\begin{aligned}
	\sum_\epsilon\lvert\epsilon_\mu^*(k)\mathcal{M}^\mu(k)\rvert^2&=\sum_{\epsilon}\epsilon_{\mu}^{*}\epsilon_{\nu}\mathcal{M}^{\mu}(k)\mathcal{M}^{\nu*}(k)\\&=|\mathcal{M}^{1}|^{2}+|\mathcal{M}^{2}|^{2}\\&=|\mathcal{M}^1|^2+|\mathcal{M}^2|^2+|\mathcal{M}^3|^2-|\mathcal{M}^0|^2\\&=-g_{\mu\nu}\mathcal{M}^{\mu}(k)\mathcal{M}^{\nu*}(k).
\end{aligned}
\]
Here, we used the Ward identity:
$$
k^\mu \mathcal{M}_\mu=k\mathcal{M}^0(k)-k\mathcal{M}^3(k)=0\Rightarrow \mathcal{M}^0=\mathcal{M}^3
$$
\textbf{MY FINAL ANSWER IS}:
\[
\boxed{
	\sum_\text{polarizations}\epsilon_\mu^*\epsilon_\nu\:\longrightarrow\:-\eta_{\mu\nu}
}
\]
\end{solution}
\begin{tcolorbox}[problem]
\textbf{6.} Compute $|{\cal M}|^2$ averaged over initial electron spin and initial photon polarizations and summed over finals, keeping $D$ symbolic up to the point where $D-2$ appears.
\end{tcolorbox}
\begin{solution}[14.6]
	\[
\begin{aligned}
	i\mathcal{M} &= \bar{u}(p')(-ie\gamma^\mu)\epsilon^*_\mu(k') \frac{i(\slashed{p} + \slashed{k} + m)}{(p+k)^2 - m^2} (-ie\gamma^\nu)\epsilon_\nu(k)u(p) \\
	&\quad + \bar{u}(p')(-ie\gamma^\nu)\epsilon_\nu(k) \frac{i(\slashed{p} - \slashed{k}' + m)}{(p-k')^2 - m^2} (-ie\gamma^\mu)\epsilon^*_\mu(k')u(p) \\[10pt]
	&= -ie^2 \epsilon^*_\mu(k') \epsilon_\nu(k) \bar{u}(p') \left[ \frac{\gamma^\mu (\slashed{p} + \slashed{k} + m) \gamma^\nu}{(p+k)^2 - m^2} + \frac{\gamma^\nu (\slashed{p} - \slashed{k}' + m) \gamma^\mu}{(p-k')^2 - m^2} \right] u(p).
\end{aligned}
\]
 Since $p^2 = m^2$ and $k^2 = 0$, the denominators of the propagators are
\[
(p+k)^2 - m^2 = 2p \cdot k \quad \text{and} \quad (p-k')^2 - m^2 = -2p \cdot k'.
\]
To simplify the numerators, we use a bit of Dirac algebra:
\begin{align*}
	(\slashed{p} + m)\gamma^\nu u(p) &= (2p^\nu - \gamma^\nu \slashed{p} + \gamma^\nu m)u(p) \\
	&= 2p^\nu u(p) - \gamma^\nu (\slashed{p} - m)u(p) \\
	&= 2p^\nu u(p).
\end{align*}
Using this trick on the numerator of each propagator, we obtain
\begin{equation*}
	i\mathcal{M} = -ie^2 \epsilon_\mu^*(k') \epsilon_\nu(k) \bar{u}(p') \left[ \frac{\gamma^\mu \slashed{k} \gamma^\nu + 2\gamma^\mu p^\nu}{2p \cdot k} + \frac{-\gamma^\nu \slashed{k}' \gamma^\mu + 2\gamma^\nu p^\mu}{-2p \cdot k'} \right] u(p). 
\end{equation*}
In the previous problem, we learned how to sum over photon polarizations. To perform the sum over fermion spins, we need the following formula:
\[
\sum_su^s(p)\bar{u}^s(p)=\slashed{p}+m,\quad\sum_sv^s(p)\bar{v}^s(p)=\slashed{p}-m
\]
\[
\begin{aligned}
	&\sum_{\text{spins}}\Big(\bar u(p')X u(p)\Big)\Big(\bar u(p)Y u(p')\Big)
	=\sum_{s,s'}\bar u_{s'}(p')X u_s(p)\,\bar u_s(p)Y u_{s'}(p')\\
	=&\sum_{s'}\bar u_{s'}(p')X\Big(\sum_s u_s(p)\bar u_s(p)\Big)Y u_{s'}(p')=\sum_{s'}\bar u_{s'}(p')X(\slashed p+m)Y u_{s'}(p')\\
	=&\operatorname{tr}\!\Big(X(\slashed p+m)Y\sum_{s'}u_{s'}(p')\bar u_{s'}(p')\Big)\\
	=&\operatorname{tr}\!\Big(X(\slashed p+m)Y(\slashed p'+m)\Big)=\operatorname{tr}\!\Big((\slashed p'+m)\,X\,(\slashed p+m)\,Y\Big).
\end{aligned}
\]
Using above formula, \textbf{MY FINAL ANSWER IS}:
\[
\begin{aligned}
	&\frac1S\sum_{\text{spins}}\sum_{\text{polarizations}}|\mathcal M|^2
	=\frac{e^4}{S}\,
	\sum_{\text{spins}}
	\Bigg[
	\Bigg(\sum_{\text{polarizations}}\epsilon_\mu^*(k')\epsilon_\rho(k')\Bigg)
	\Bigg(\sum_{\text{polarizations}}\epsilon_\nu(k)\epsilon_\sigma^*(k)\Bigg)
	\\
	&\qquad\qquad\qquad\qquad\qquad\qquad\times
	\Bigg(\bar u(p')\Bigg[
	\frac{\gamma^\mu\slashed k\,\gamma^\nu+2\gamma^\mu p^\nu}{2p\cdot k}
	+\frac{\gamma^\nu\slashed k'\,\gamma^\mu-2\gamma^\nu p^\mu}{2p\cdot k'}
	\Bigg]u(p)\Bigg)
	\\
	&\qquad\qquad\qquad\qquad\qquad\qquad\times
	\Bigg(\bar u(p')\Bigg[
	\frac{\gamma^\rho\slashed k\,\gamma^\sigma+2\gamma^\rho p^\sigma}{2p\cdot k}
	+\frac{\gamma^\sigma\slashed k'\,\gamma^\rho-2\gamma^\sigma p^\rho}{2p\cdot k'}
	\Bigg]u(p)\Bigg)^*
	\Bigg]
	\\
	&=\frac{e^4}{S}\,
	\eta_{\mu\rho}\,\eta_{\nu\sigma}\,
	\sum_{\text{spins}}
	\Bigg[
	\bar u(p')\Bigg[
	\frac{\gamma^\mu\slashed k\,\gamma^\nu+2\gamma^\mu p^\nu}{2p\cdot k}
	+\frac{\gamma^\nu\slashed k'\,\gamma^\mu-2\gamma^\nu p^\mu}{2p\cdot k'}
	\Bigg]u(p)\,
	\bar u(p)\,
	\Bigg[
	\frac{\gamma^\sigma\slashed k\,\gamma^\rho+2\gamma^\rho p^\sigma}{2p\cdot k}
	+\frac{\gamma^\rho\slashed k'\,\gamma^\sigma-2\gamma^\sigma p^\rho}{2p\cdot k'}
	\Bigg]u(p')
	\Bigg]
	\\
	&=\frac{e^4}{S}\,
	\eta_{\mu\rho}\,\eta_{\nu\sigma}\,
	\operatorname{tr}\Bigg\{
	(\slashed p'+m)\Bigg[
	\frac{\gamma^\mu\slashed k\,\gamma^\nu+2\gamma^\mu p^\nu}{2p\cdot k}
	+\frac{\gamma^\nu\slashed k'\,\gamma^\mu-2\gamma^\nu p^\mu}{2p\cdot k'}
	\Bigg]
	(\slashed p+m)\Bigg[
	\frac{\gamma^\sigma\slashed k\,\gamma^\rho+2\gamma^\rho p^\sigma}{2p\cdot k}
	+\frac{\gamma^\rho\slashed k'\,\gamma^\sigma-2\gamma^\sigma p^\rho}{2p\cdot k'}
	\Bigg]
	\Bigg\}
\end{aligned}
\]
In $D$ dimensions, a Dirac spinor has $2^{\lfloor D/2\rfloor}$ degrees of freedom. Since the electron is a massive fermion and satisfies the Dirac equation, the number of on-shell degrees of freedom is $2^{\lfloor D/2\rfloor-1}$. The photon has $D-2$ on-shell degrees of freedom, so in the above expression
\[
S=2^{\lfloor D/2\rfloor-1}\cdot(D-2).
\]

\end{solution}

\begin{tcolorbox}[problem]
\textbf{7.} Derive the relation between Mandelstam variables $s,t,u$ for Compton scattering with a massive electron, $s+t+u=2m^2$.
\end{tcolorbox}
\begin{solution}[14.7]
	The convention for momentum is $\gamma(k) + e^-(p) \to \gamma(k') + e^-(p')$. Using the definition of Mandelstam variables and the on-shell conditions ($p^2=p'^2=m^2$ and $k^2=k'^2=0$), we have:
	\begin{align*}
		s &= (p + k)^2 = p^2 + k^2 + 2p \cdot k = m^2 + 2p \cdot k, \\
		t &= (k - k')^2 = k^2 + k'^2 - 2k \cdot k' = -2k \cdot k', \\
		u &= (p - k')^2 = p^2 + k'^2 - 2p \cdot k' = m^2 - 2p \cdot k'.
	\end{align*}
	Summing these three equations yields:
	\begin{equation*}
		s + t + u = 2m^2 + 2(p \cdot k - k \cdot k' - p \cdot k').
	\end{equation*}
	From the conservation of four-momentum, $p + k = p' + k'$, which can be rearranged as $p' = p + k - k'$. Squaring both sides leads to:
	\begin{equation*}
		p'^2 = (p + k - k')^2 = p^2 + k^2 + k'^2 + 2(p \cdot k - p \cdot k' - k \cdot k').
	\end{equation*}
	Applying the on-shell conditions again, the equation simplifies to:
	\begin{equation*}
		m^2 = m^2 + 2(p \cdot k - p \cdot k' - k \cdot k') \implies 2(p \cdot k - p \cdot k' - k \cdot k') = 0.
	\end{equation*}
	Substituting this result back into the expression for the sum, we arrive at the required relation:
	\begin{equation*}
		s + t + u = 2m^2.
	\end{equation*}
\end{solution}

\begin{tcolorbox}[problem]
\textbf{8.} Work out the $m\to 0$ limit of the squared amplitude and identify which kinematic regions generate collinear enhancement.
\end{tcolorbox}
\begin{solution}[14.8]
	Using the result in problem 9, we have:
	\[
	\mathcal{A}\xrightarrow{m\to 0}2e^4\left[\frac{p\cdot k^{\prime}}{p\cdot k}+\frac{p\cdot k}{p\cdot k^{\prime}}\right]
	\]
	In massless limit, i.e. high energy limit:
	\[
	\begin{aligned}
			p\cdot k
		&= p^\mu k_\mu
		= p^0k^0-\vec p\cdot\vec k
		= E_pE_k-|\vec p||\vec k|\cos\theta_{pk}
		= E_pE_k-|\vec p|E_k\cos\theta_{pk}
		= E_k\bigl(E_p-|\vec p|\cos\theta_{pk}\bigr)\\
		&\xrightarrow[m\to0]{\,|\vec p|\to E_p\,}
		E_k\bigl(E_p-E_p\cos\theta_{pk}\bigr)
		= E_pE_k(1-\cos\theta_{pk})
	\end{aligned}
	\]
	Similarly,
	\[
	p\cdot k' \xrightarrow[m\to0]{} E_pE_{k'}(1-\cos\theta_{pk'})
	\]
	Therefore, $p\cdot k\sim 0$ means $\theta_{pk}\sim 0$, i.e. the initial electron parallels to the final photon ($k\parallel p$). so the collinear enhancement arises from the regions \(p\cdot k\to0\) (i.e.\ \(k\parallel p\)) and/or \(p\cdot k'\to0\) (i.e.\ \(k'\parallel p\)).
	
\end{solution}
\begin{tcolorbox}[problem]
\textbf{9.} Specialize your general $D$-dimensional expression to $D=4$ and recover the standard Klein--Nishina structure (you can quote a known final form, but you must show how your polarization counting matches).
\end{tcolorbox}
\begin{solution}[14.9]
	In four dimention, $S=2^{2-1}\cdot(4-2)=4$ in problem 4. There are 4 terms in the amplitude:
	$$
	\mathcal{A}=\frac{e^4}{4}\left[\frac{\mathbf{I}}{(2p\cdot k)^2}+\frac{\mathbf{II}}{(2p\cdot k)(2p\cdot k^{\prime})}+\frac{\mathbf{III}}{(2p\cdot k^{\prime})(2p\cdot k)}+\frac{\mathbf{IV}}{(2p\cdot k^{\prime})^2}\right]
	$$
	The first term is:
	\[
	\mathbf{I}=\mathrm{tr}\left[(p^{\prime}+m)(\gamma^{\mu}\not k\gamma^{\nu}+2\gamma^{\mu}p^{\nu})(\not p+m)(\gamma_{\nu}\not k\gamma_{\mu}+2\gamma_{\mu}p_{\nu})\right]
	\]
	The most difficult part is:
\[
\begin{aligned}
	\operatorname{tr}\!\left[\slashed{p}'\,\gamma^\mu\,\slashed{k}\,\gamma^\nu\,\slashed{p}\,\gamma_\nu\,\slashed{k}\,\gamma_\mu\right]
	&= \operatorname{tr}\!\left[(-2\slashed{p}')\,\slashed{k}\,(-2\slashed{p})\,\slashed{k}\right] \\
	&= \operatorname{tr}\!\left[4\,\slashed{p}'\,\slashed{k}\,\left(2\,p\!\cdot\! k-\slashed{k}\slashed{p}\right)\right] \\
	&= 8\,p\!\cdot\! k\ \operatorname{tr}\!\left[\slashed{p}'\,\slashed{k}\right] \\
	&= 32\,(p\!\cdot\! k)\,(p'\!\cdot\! k)\,.
\end{aligned}
\]
	Therefore we obtain:
	\[
	\mathbf{I}=16\left(4m^4-2m^2p\cdotp^{\prime}+4m^2p\cdot k-2m^2p^{\prime}\cdot k+2(p\cdot k)(p^{\prime}\cdot k)\right)
	\]
	we can use Mandelstam variables to rewrite it as:
	\[
	\mathbf{I}=16{\left(2m^4+m^2(s-m^2)-\frac{1}{2}(s-m^2)(u-m^2)\right)}
	\]
	Sending $k\leftrightarrow-k^\prime$, we can immediately write
	$$
	\mathbf{IV}=16\big(2m^4+m^2(u-m^2)-\frac12(s-m^2)(u-m^2)\big).
	$$
	Similarly, the answer of $\mathbf{II}$ and $\mathbf{III}$ is
	$$\mathbf{II}=\mathbf{III}=-8\big(4m^4+m^2(s-m^2)+m^2(u-m^2)\big).$$
	Then we have:
	\[
	\mathcal{A}=2e^4\left[\frac{p\cdot k^{\prime}}{p\cdot k}+\frac{p\cdot k}{p\cdot k^{\prime}}+2m^2\left(\frac{1}{p\cdot k}-\frac{1}{p\cdot k^{\prime}}\right)+m^4\left(\frac{1}{p\cdot k}-\frac{1}{p\cdot k^{\prime}}\right)^2\right]
	\]
	To derive the Klein-Nishina formula, we need to choose the following “lab” frame:
	\[
			k=(\omega,\omega\hat{z}),\quad p=(m,\mathbf{0}),\quad
			k^{\prime}=(\omega^{\prime},\omega^{\prime}\sin\theta,0,\omega^{\prime}\cos\theta),\quad p^{\prime}=(E^{\prime},\mathbf{p}^{\prime})
	\]
	The phase space integral in this frame is
	\[
	\begin{aligned}\int d\Pi_{2}&=\int\frac{d^3k^{\prime}}{(2\pi)^3}\frac{1}{2\omega^{\prime}}\frac{d^3p^{\prime}}{(2\pi)^3}\frac{1}{2E^{\prime}}\left(2\pi\right)^4\delta^{(4)}(k^{\prime}+p^{\prime}-k-p)\\&=\int\frac{(\omega^{\prime})^2d\omega^{\prime}d\Omega}{(2\pi)^3}\frac{1}{4\omega^{\prime}E^{\prime}}\\&\times2\pi\delta(\omega^{\prime}+\sqrt{m^2+\omega^2+(\omega^{\prime})^2-2\omega\omega^{\prime}\cos\theta}-\omega-m)\\&=\int\frac{d\left.\cos\theta\right.}{2\pi}\frac{\omega^{\prime}}{4E^{\prime}}\frac{1}{\left|1+\frac{\omega^{\prime}-\omega\cos\theta}{E^{\prime}}\right|}\\&=\frac{1}{8\pi}\int d\left.\cos\theta\right.\frac{\omega^{\prime}}{m+\omega(1-\cos\theta)}\\&=\frac{1}{8\pi}\int d\cos\theta\frac{(\omega^{\prime})^{2}}{\omega m}.\end{aligned}
	\]
	Plugging everything in the Golden formula:
	\[
	\begin{aligned}
		d\sigma&=\frac{1}{2E_{\mathcal{A}}2E_{\mathcal{B}}\left|v_{\mathcal{A}}-v_{\mathcal{B}}\right|}\left(\prod_{f}\frac{d^{3}p_{f}}{(2\pi)^{3}}\frac{1}{2E_{f}}\right)\\&\times\left|\mathcal{M}(p_{\mathcal{A}},p_{\mathcal{B}}\to\{p_{f}\})\right|^{2}(2\pi)^{4}\delta^{(4)}(p_{\mathcal{A}}+p_{\mathcal{B}}-\sum p_{f}).
		\end{aligned}
	\]
	Here, we need to set $\left|v_{\mathcal{A}}-v_{\mathcal{B}}\right|=1$. \textbf{MY FINAL ANSWER IS}:
	\[
	\boxed{
	\frac{d\sigma}{d\cos\theta}=\frac{1}{2\omega}\frac{1}{2m}\cdot\frac{1}{8\pi}\frac{(\omega^{\prime})^2}{\omega m}\cdot\mathcal{A}=\frac{\pi\alpha^2}{m^2}\left(\frac{\omega^{\prime}}{\omega}\right)^2\left[\frac{\omega^{\prime}}{\omega}+\frac{\omega}{\omega^{\prime}}-\sin^2\theta\right]}
	\]
\end{solution}

\begin{tcolorbox}[problem]
\textbf{10.} Consider $D=3$ ($2+1$ dimensions), how does the absence of a second transverse polarization change the spin/polarization averages in $|{\cal M}|^2$? Write the modified averaging factor and explain qualitatively how the cross section changes.
\end{tcolorbox}
\begin{solution}[14.10]
	The modified averaging factor is:
	\[
	\boxed{
	S=\frac12
	}
	\]
\end{solution}

\end{document}

	\clearpage
	\documentclass[11pt]{article}

\usepackage[margin=1.6cm]{geometry}
\usepackage{amsmath,amssymb,amsthm}
\usepackage[hidelinks]{hyperref}
\usepackage[most]{tcolorbox}
\usepackage[]{cancel}
\usepackage{tikz-feynman}
\usepackage{slashed}
\tikzfeynmanset{compat=1.1.0}
\newenvironment{solution}[1][]{\begin{proof}[\textbf{Solution}]}{\end{proof}}
% --- Rounded box style for each problem (whole statement) ---
\tcbset{
    problem/.style={
        enhanced,
        boxrule=0.7pt,
        arc=3mm,
        left=3mm,right=3mm,top=2mm,bottom=2mm,
        colback=white,
        colframe=black,
        before skip=10pt,
        after skip=6pt,
    }
}

\begin{document}
\begin{center}
    {\LARGE \textbf{Solution Manual of Problem Set 15} \par}
    \vspace{1em}
    {\large Bufan Zheng\\ \href{mailto:bufan@g.ecc.u-tokyo.ac.jp}{\texttt{bufan@g.ecc.u-tokyo.ac.jp}} \par}
\end{center}


\begin{tcolorbox}[problem]
\textbf{1.} Define ``soft divergence'' and ``collinear divergence'' in terms of regions of loop or phase-space integration.
\end{tcolorbox}
\begin{solution}[15.1]
	Soft singularities occur when the momentum of a massless particle goes to zero. Feynman diagrams of the form
	\[
	\begin{tikzpicture}[baseline=(current bounding box.center)]
		\begin{feynman}
			% 1. 定义中心标准 blob
			% 可以通过 minimum size 稍微调整它的大小
			\vertex [blob, minimum size=1cm] (c) {};
			
			% 2. 定义周围的顶点位置（手动指定以模拟原图的随意感）
			% 左侧入射
			\vertex [above left=2cm of c] (i1);
			\vertex [left=1.8cm of c, yshift=0.5cm] (i2);
			\vertex [below left=2.2cm of c, yshift=1cm] (i3);
			\vertex [below left=1cm of c,yshift=-1.2cm, rotate=-20] (i4);
			
			% 右侧出射
			% 关键的带动量 k 的光子顶点位置
			\vertex [above right=2.5cm of c] (o_k); 
			\vertex [above right=1.5cm of c] (o_f); 
			\vertex [right=2.5cm of c,yshift=0.5cm] (o_f2); 
			% 其他出射粒子
			\vertex [below right=2cm of c] (o2);
			\vertex [right=1.8cm of c, yshift=-0.5cm] (o3);
			
			% 3. 绘制连接线 (使用 \diagram* 手动模式)
			\diagram* {
				% 入射粒子 (混合直线和波浪线)
				(i1) -- [fermion] (c),
				(i2) -- [photon] (c),
				(i3) -- [photon] (c),
				(i4) -- [fermion] (c),
				
				% 关键：带有动量 k 的出射光子
				% momentum 选项会自动添加箭头和标签
				(c) -- [fermion, momentum=$\ell$] (o_f),
				(o_f) -- [fermion,momentum=$p$](o_k),
				(o_f) -- [photon, momentum'=$k$](o_f2),
				
				% 其他出射粒子
				(c) -- [fermion] (o2),
				(c) -- [fermion] (o3),
			};
		\end{feynman}
	\end{tikzpicture}
	\]
	develop a singularity in this limit because the propagator, goes on-shell, i.e. $\ell=k+p\to p$.
	
	Another type of amplitude singularity arises when two external momenta become collinear. Diagrams in which the two external legs attach to a common 3-vertex are then divergent because the internal propagator goes on-shell. Still using the diagram above as an example, we assume the electron is massless. When the emitted photon momentum is parallel to the outgoing electron momentum, namely \(k \propto p\), we clearly have \(\ell \propto k \propto p\). Since \(k^2 = 0\), it follows that the propagator satisfies \(\ell^2 = 0\), meaning the internal propagator goes on shell, which leads to a divergence.
\end{solution}
\begin{tcolorbox}[problem]
\textbf{2.} Show by power counting in $D$ dimensions that the phase space measure for emission of a soft massless boson behaves like $d^Dk\,\delta(k^2)\theta(k^0)\sim k^{D-3}dk$. For which $D$ does $$\int_0 dk\,\frac{k^{D-3}}{k}$$ diverge? (soft factor $\sim 1/k$)
\end{tcolorbox}
\begin{solution}[15.2]
	The phase measure for massless boson behaves like:
	\[
	\begin{aligned}
	d^Dk\,\delta(k^2)\,\theta(k^0)
	&= dk^0\,d^{D-1}\mathbf{k}\,\delta\!\big((k^0)^2-\mathbf{k}^2\big)\,\theta(k^0)
	= dk^0\,d^{D-1}\mathbf{k}\,\frac{\delta(k^0-|\mathbf{k}|)}{2|\mathbf{k}|}\\
	&= \frac{d^{D-1}\mathbf{k}}{2|\mathbf{k}|}
	= \frac{k^{D-2}dk\,d\Omega_{D-2}}{2k}
	\sim k^{D-3}\,dk .
	\end{aligned}
	\]
	Next, consider the soft limit, we encounter the following integral:
	\[
	\int_0dk\frac{k^{D-3}}{k}\sim k^{D-3}
	\]
	It is clear the divergence appear when $D-3\leq 0$. So \textbf{MY FINAL ANSWER IS}:
	\[
	\boxed{
	D\leq 3
	}
	\]
\end{solution}
\begin{tcolorbox}[problem]
\textbf{3.} For massless charged particles in $D=4$, explain why collinear divergences appear but are absent if the charged particle has a nonzero mass.
\end{tcolorbox}
\begin{solution}[15.3]
	From the discussion in Problem~1, we can see that regardless of whether the electron is massless, when the momentum of the radiated photon approaches zero, the internal electron propagator goes on shell, i.e.\ \(\ell^2 = m^2\). However, when we consider the collinear limit, we can only obtain \(\ell^2 \propto k^2 = 0\). Therefore, only when the electron is massless can we say that the propagator goes on shell and thus produces a divergence. If the electron is massive, the denominator of the propagator still receives a nonzero contribution of \(m^2\).
\end{solution}
\begin{tcolorbox}[problem]
\textbf{4.} Derive the eikonal factor for soft photon emission off an external leg, $\sim e\,\frac{p\cdot\epsilon}{p\cdot k}$. Explain why it is universal.
\end{tcolorbox}
\begin{solution}[15.4]
	Consider the following diagram:
		\[
	\begin{tikzpicture}[baseline=(current bounding box.center)]
		\begin{feynman}
			% 1. 定义中心标准 blob
			% 可以通过 minimum size 稍微调整它的大小
			\vertex [blob, minimum size=1cm] (c) {};
			
			% 2. 定义周围的顶点位置（手动指定以模拟原图的随意感）
			% 左侧入射
			\vertex [above left=2cm of c] (i1);
			\vertex [left=1.8cm of c, yshift=0.5cm] (i2);
			\vertex [below left=2.2cm of c, yshift=1cm] (i3);
			\vertex [below left=1cm of c,yshift=-1.2cm, rotate=-20] (i4);
			
			% 右侧出射
			% 关键的带动量 k 的光子顶点位置
			\vertex [above right=2.5cm of c] (o_k); 
			\vertex [above right=1.5cm of c] (o_f); 
			\vertex [right=2.5cm of c,yshift=0.5cm] (o_f2); 
			% 其他出射粒子
			\vertex [below right=2cm of c] (o2);
			\vertex [right=1.8cm of c, yshift=-0.5cm] (o3);
			
			% 3. 绘制连接线 (使用 \diagram* 手动模式)
			\diagram* {
				% 入射粒子 (混合直线和波浪线)
				(i1) -- [fermion] (c),
				(i2) -- [photon] (c),
				(i3) -- [photon] (c),
				(i4) -- [fermion] (c),
				
				% 关键：带有动量 k 的出射光子
				% momentum 选项会自动添加箭头和标签
				(c) -- [fermion, momentum=$\ell$] (o_f),
				(o_f) -- [fermion,momentum=$p$](o_k),
				(o_f) -- [photon, momentum'=$k$](o_f2),
				
				% 其他出射粒子
				(c) -- [fermion] (o2),
				(c) -- [fermion] (o3),
			};
		\end{feynman}
	\end{tikzpicture}
	\quad=\quad
	\begin{tikzpicture}[baseline=(current bounding box.center)]
		\begin{feynman}
			% 1. 定义中心标准 blob
			% 可以通过 minimum size 稍微调整它的大小
			\vertex [blob, minimum size=2cm] (c) {};
			
			% 2. 定义周围的顶点位置（手动指定以模拟原图的随意感）
			
			% 右侧出射
			% 关键的带动量 k 的光子顶点位置
			\vertex [above right=3.5cm of c] (o_k); 
			\vertex [above right=2cm of c] (o_f); 
			\vertex [right=2.5cm of c,yshift=1cm] (o_f2); 
			
			% 3. 绘制连接线 (使用 \diagram* 手动模式)
			\diagram* {
				
				% 关键：带有动量 k 的出射光子
				% momentum 选项会自动添加箭头和标签
				(c) -- [fermion, momentum=$\ell$] (o_f),
				(o_f) -- [fermion,momentum=$p$](o_k),
				(o_f) -- [photon, momentum'=$k$](o_f2),
			};
		\end{feynman}
	\end{tikzpicture}
	\]
	For scalar QED, we have,
	\[
	\begin{aligned}
	&\mathcal{M}_{\text{soft}}=\epsilon_{\mu}\cdot(i e (2p^{\mu}+k^{\mu}))\cdot\frac{-i}{(p+k)^2-m^2+i\varepsilon}\cdot\mathcal{M}_{\text{hard}}\\
	&\xrightarrow{k\to 0}\epsilon_\mu e(2p^\mu)\frac{1}{2p\cdot k+i\varepsilon}\left.\mathcal{M}_{\mathrm{hard}}\xrightarrow{\varepsilon\to 0} e\frac{\epsilon\cdot p}{p\cdot k}\right.\mathcal{M}_{\mathrm{hard}}
	\end{aligned}
	\]
	For fermion QED, we have,
	\[
	\mathcal{M}_{\text{soft}}=\bar{u}(p)\cdot\epsilon_{\mu}\cdot(i e \gamma^{\mu})\cdot\frac{-i(\slashed{p}+\slashed{k}+m)}{(p+k)^2-m^2+i\varepsilon}\cdot\hat{\mathcal{M}}_{\text{hard}}
	\]
	Note here $\bar{u}(p+k)\cdot \hat{\mathcal{M}}_{\text{hard}}={\mathcal{M}}_{\text{hard}}$. Use $\bar u(p)(\slashed p-m)=0$ and the identity
	\[
	\slashed{\epsilon}\slashed p = 2\,\epsilon\!\cdot\! p-\slashed p\slashed{\epsilon}
	\]
	to get
	\[
	\bar u(p)\slashed{\epsilon}(\slashed p+m)
	=\bar u(p)\bigl(\slashed{\epsilon}\slashed p+m\slashed{\epsilon}\bigr)
	=2\,\epsilon\!\cdot\! p\,\bar u(p).
	\]
	Therefore
	\[
	\bar u(p)\slashed{\epsilon}(\slashed p+\slashed k+m)
	=2\,\epsilon\!\cdot\! p\,\bar u(p)+\bar u(p)\slashed{\epsilon}\slashed k,
	\]
	plug in the relation between $\mathcal{M}_{\text{soft}}$ and $\mathcal{M}_{\text{hard}}$:
	\[
	\mathcal{M}_{\text{soft}}=e\frac{2\epsilon\cdot p }{2p\cdot k +i\varepsilon}\bar{u}(p)\hat{\mathcal{M}}_{\text{hard}}+\mathcal{O}(k)\sim e\frac{2\epsilon\cdot p }{2p\cdot k }\bar{u}(p+k)\hat{\mathcal{M}}_{\text{hard}}\sim e\frac{\epsilon\cdot p}{p\cdot k}\cdot\mathcal{M}_{\text{hard}}
	\]
	So, \textbf{MY FINAL ANSWER IS}:
	\[
	\boxed{
		\mathcal{M}_{\text{soft}}\sim e\frac{\epsilon\cdot p}{p\cdot k}\cdot\mathcal{M}_{\text{hard}}
	}
	\]
	
\end{solution}
\begin{tcolorbox}[problem]
\textbf{5.} Show how the soft emission probability factorizes and exponentiates (Bloch--Nordsieck / Yennie--Frautschi--Suura idea) at leading soft order.
\end{tcolorbox}

\begin{tcolorbox}[problem]
\textbf{6.} In $D=4$, show schematically how the virtual soft divergence cancels the real soft divergence in an inclusive cross section $\sigma_{\text{incl}}$ (you don't need full constants, demonstrate matching singular structures).
\end{tcolorbox}

\begin{tcolorbox}[problem]
\textbf{7.} Define ``infrared safe observable'' and give two examples in QED and two in QCD (briefly).
\end{tcolorbox}

\begin{tcolorbox}[problem]
\textbf{8.} Explain what changes when you go from fully inclusive to ``semi-inclusive'' observables (e.g.\ restricting photon energy or restricting angles). Which restrictions reintroduce logarithms?
\end{tcolorbox}

\begin{tcolorbox}[problem]
\textbf{9.} Sudakov logs, show (by region analysis) why double logs $\alpha\ln^2(Q^2/m^2)$ arise in $D=4$ from overlapping soft+collinear regions.
\end{tcolorbox}

\begin{tcolorbox}[problem]
\textbf{10.} In $D=4$, specify exactly which combination is IR finite order-by-order for a process with charged particles: (a) fixed-order exclusive $e^-e^-\to e^-e^-$ with no photons, (b) inclusive over any number of photons below an energy resolution $\Delta E$, (c) ``dressed electron'' definition where photons within a small cone are recombined with the electron. For each, state whether it is IR finite and why (one paragraph total).
\end{tcolorbox}


\end{document}

	\clearpage
	\documentclass[11pt]{article}

\usepackage[margin=1.6cm]{geometry}
\usepackage{amsmath,amssymb,amsthm}
\usepackage[hidelinks]{hyperref}
\usepackage[most]{tcolorbox}
\usepackage[]{cancel}
\usepackage{tikz-feynman}
\tikzfeynmanset{compat=1.1.0}
\newenvironment{solution}[1][]{\begin{proof}[\textbf{Solution}]}{\end{proof}}
% --- Rounded box style for each problem (whole statement) ---
\tcbset{
    problem/.style={
        enhanced,
        boxrule=0.7pt,
        arc=3mm,
        left=3mm,right=3mm,top=2mm,bottom=2mm,
        colback=white,
        colframe=black,
        before skip=10pt,
        after skip=6pt,
    }
}


\begin{document}
\begin{center}
    {\LARGE \textbf{Solution Manual of Problem Set 16} \par}
    \vspace{1em}
    {\large Bufan Zheng\\ \href{mailto:bufan@g.ecc.u-tokyo.ac.jp}{\texttt{bufan@g.ecc.u-tokyo.ac.jp}} \par}
\end{center}

\begin{tcolorbox}[problem]
\textbf{1.} Write the most general Lorentz-covariant decomposition of the on-shell QED vertex $\Gamma^\mu(p',p)$ in $D=4$ (Dirac + Pauli form factors).
\end{tcolorbox}
\begin{solution}[16.1]
	\textbf{MY FINAL ANSWER IS}
	\[
\boxed{\bar{u}(p^{\prime})\operatorname{\Gamma}^\mu(p^{\prime},p)\operatorname{u}(p)=\bar{u}(p^{\prime})\left[F_1(q^2)\gamma^\mu+i\operatorname{}\frac{F_2(q^2)}{2m}\operatorname{\sigma}^{\mu\nu}q_\nu\right]u(p)}
	\]
	Where, $\sigma^{\mu\nu}\equiv\frac{i}{2}[\gamma^\mu,\gamma^\nu]$.
\end{solution}

\begin{tcolorbox}[problem]
\textbf{2.} Explain what changes in general $D$, what replaces $\sigma^{\mu\nu}$ as an independent structure, and what assumptions (on-shell, parity, time reversal) constrain the decomposition.
\end{tcolorbox}

\begin{tcolorbox}[problem]
\textbf{3.} Define $a_e$ (``anomalous magnetic moment'') in a way that makes sense as $D\to 4$ in dimensional regularization, and explain why the physical observable is extracted from $F_2(0)$ after renormalization.
\end{tcolorbox}

\begin{tcolorbox}[problem]
\textbf{4.} Compute the superficial degree of divergence of the one-loop vertex correction in $D$ dimensions.
\end{tcolorbox}
\begin{solution}[16.4]
	The amplitude corresponding to a typical Feynman diagram has the following structure. 
	\[
\int\frac{d^Dk_1d^Dk_2\cdots d^Dk_L}{(k_i-m)\cdots(k_j^2)\cdots(k_n^2)}.
	\]
	The superficial degree of divergence $\mathfrak{D}$ can be calculated using the formula below. 
	\[
	\begin{aligned}\mathfrak{D}&\equiv(\text{power of }k\text{ in numerator})-(\text{power of }k\text{ in denominator})\\&=DL-P_e-2P_\gamma.\end{aligned}
	\]
	For the one-loop vertex correction in QED, $P_e=2$,$P_\gamma=1$,$L=1$ , so we obtain
	\[
	\boxed{
	\mathfrak{D}= D-2-2=D-4}
	\]
\end{solution}

\begin{tcolorbox}[problem]
\textbf{5.} Show that the one-loop vertex correction satisfies the Ward--Takahashi identity relating it to the electron self-energy.
\end{tcolorbox}

\begin{tcolorbox}[problem]
\textbf{6.} In covariant $\xi$ gauge, identify which terms in the photon propagator can potentially contribute to $F_2(0)$, and argue why the final $a_e$ is gauge independent.
\end{tcolorbox}

\begin{tcolorbox}[problem]
\textbf{7.} In $D=4-2\epsilon$, show how UV divergences appear as $1/\epsilon$ poles in the vertex and how they are removed by renormalization.
\end{tcolorbox}

\begin{tcolorbox}[problem]
\textbf{8.} Explain why $F_2(0)$ is finite in renormalized QED at one loop in $D=4$ even though individual integrals contain UV divergences.
\end{tcolorbox}

\begin{tcolorbox}[problem]
\textbf{9.} Derive (or outline derivation of) the classic one-loop result $a_e=\alpha/(2\pi)$ in $D=4$, identifying where the ``$1/2$'' comes from (trace algebra + Feynman parameters).
\end{tcolorbox}

\begin{tcolorbox}[problem]
\textbf{10.} Discuss dimension dependence, for which $D$ would you expect the one-loop contribution to $a_e$ to be power-law UV divergent instead of logarithmic, based on operator dimension/power counting?
\end{tcolorbox}


\end{document}

	\clearpage
	\documentclass[11pt]{article}
\usepackage[margin=1.6cm]{geometry}
\usepackage{amsmath,amssymb,amsthm}
\allowdisplaybreaks[4]
\usepackage[hidelinks]{hyperref}
\usepackage[most]{tcolorbox}
\usepackage[]{cancel}
\usepackage{tikz-feynman}
\tikzfeynmanset{compat=1.1.0}
\newenvironment{solution}[1][]{\begin{proof}[\textbf{Solution}]}{\end{proof}}
% --- Rounded box style for each problem (whole statement) ---
\tcbset{
    problem/.style={
        enhanced,
        boxrule=0.7pt,
        arc=3mm,
        left=3mm,right=3mm,top=2mm,bottom=2mm,
        colback=white,
        colframe=black,
        before skip=10pt,
        after skip=6pt,
    }
}


\begin{document}
\begin{center}
    {\LARGE \textbf{Solution Manual of Problem Set 17} \par}
    \vspace{1em}
    {\large Bufan Zheng\\ \href{mailto:bufan@g.ecc.u-tokyo.ac.jp}{\texttt{bufan@g.ecc.u-tokyo.ac.jp}} \par}
\end{center}


\begin{tcolorbox}[problem]
\textbf{1.} For a massless scalar on an interval of length $L$ with Dirichlet boundary conditions in $D=2$, write the normal mode frequencies $\omega_n$.
\end{tcolorbox}
\begin{solution}[17.1]
		We consider a massless scalar field $\phi(t,x)$ on the interval $x\in[0,L]$ in $D=2$,
		satisfying the wave equation
		\[
		(\partial_t^2-\partial_x^2)\phi(t,x)=0,
		\]
		with Dirichlet boundary conditions
		\[
		\phi(t,0)=0,\qquad \phi(t,L)=0.
		\]
		Separating variables $\phi(t,x)=e^{-i\omega t}f(x)$ gives
		\[
		f''(x)+\omega^2 f(x)=0,\qquad f(0)=f(L)=0,
		\]
		whose normalizable solutions are
		\[
		f_n(x)=\sin\!\left(\frac{n\pi x}{L}\right),\qquad n=1,2,3,\ldots
		\]
		Thus the allowed wave numbers and normal-mode frequencies are
		\[
		k_n=\frac{n\pi}{L},\qquad \omega_n=|k_n|=\frac{n\pi}{L},\qquad n=1,2,3,\ldots
		\]
		So, \textbf{MY FINAL ANSWER IS}:
		\[
		\boxed{
		\omega_n=|k_n|=\frac{n\pi}{L},\quad n=1,2,3,\ldots
		}
		\]
\end{solution}
\begin{tcolorbox}[problem]
\textbf{2.} Formally write the vacuum energy $E_0(L)=\frac12\sum_n\omega_n$ and explain why it diverges.
\end{tcolorbox}
\begin{solution}[17.2]
	Using the frequencies found in Problem 1, we obtain the formal expression
	\[
	E_0(L)=\frac{1}{2}\sum_{n=1}^{\infty}\frac{n\pi}{L}
	=\frac{\pi}{2L}\sum_{n=1}^{\infty} n.
	\]
	This diverges because the high-frequency (large-$n$) modes contribute without bound:
	as $n\to\infty$, $\omega_n\sim n$, so the partial sums grow like
	\[
	\sum_{n=1}^{N}\omega_n \sim \frac{\pi}{L}\sum_{n=1}^{N} n
	= \frac{\pi}{L}\,\frac{N(N+1)}{2}\xrightarrow[N\to\infty]{}\infty.
	\]
	Equivalently, $\sum_{n=1}^{\infty} n$ is a divergent series, reflecting the ultraviolet divergence of the vacuum energy.
\end{solution}
\begin{tcolorbox}[problem]
\textbf{3.} Use zeta-function regularization in $D=2$ to compute the finite part of $E_0(L)$ and show it scales like $-1/L$.
\end{tcolorbox}
\begin{solution}[17.3]
	The zeta-function regularization is:
	\[
	1+2+3+\cdots=\zeta(-1)=-\frac{1}{12}
	\]
	plug it into the result of last problem:
	\[
	E_0(L)=\frac{\pi}{2L}\sum_{n=1}^{\infty}n=\frac{\pi}{2L}\cdot-\frac{1}{12}+\infty=-\frac{\pi}{24L}+\infty
	\]
	It is clear it scales like $-\frac{1}{L}$.
\end{solution}
\begin{tcolorbox}[problem]
\textbf{4.} Explain precisely what ``subtracting local counterterms'' means here, identify the terms in $E_0$ that are extensive in $L$ and how they are absorbed into local terms (cosmological constant / boundary tensions if relevant).
\end{tcolorbox}
\begin{solution}[17.4]
	As a first step, we replace $E_0$ by the regularized quantity:
	\[
	E_0\left(L;\alpha\right)=\frac{\pi}{2L}\sum_{n=1}^\infty n\exp\left[-\frac{n\alpha}{L}\right]
	\]
	where $\alpha$ is the cutoff parameter. It is easy to see that this series converges for $\alpha>0$, while the original divergent expression is recovered in the limit $\alpha\to0$. A straightforward calculation gives
	\[
	\begin{aligned}
	E_0\left(L;\alpha\right)&=-\frac{\pi}{2}\frac{\partial}{\partial\alpha}\sum_{n=1}^\infty\exp\left[-\frac{n\alpha}{L}\right]=\frac{\pi}{2L}\frac{\exp\left(-\frac{\alpha}{L}\right)}{\left[1-\exp\left(-\frac{\alpha}{L}\right)\right]^2}\\
	&=\frac{\pi}{8L}\frac{1}{\sinh^2\frac{\alpha}{2L}}=\frac{\pi L}{2\alpha^2}-\frac{\pi}{24L}+\frac{1}{L}O\left(\frac{\alpha^2}{L^2}\right)
	\end{aligned}
	\]
	As $\alpha\to0$, the first term here diverges as $\alpha^{-2}$, the second term remains finite and further terms vanish. The crucial fact is that the singular term does not depend on the distance $L$ between the plates and \textbf{can be interpreted as the energy density of the zero-point fluctuations in Minkowski spacetime without boundaries.} 
	\[
	S_{\rm ct}=-\int dt\int_{0}^{L}dx\;\delta\Lambda
	\]
	It is clear this is a local counter term, so the divergent part can be absorbed by local counter term \textbf{due to the cosmological constant.}
\end{solution}
\begin{tcolorbox}[problem]
\textbf{5.} Show that the remaining finite part depends nonlocally on $L$ and cannot be cancelled by any counterterm that is an integral of local densities.
\end{tcolorbox}
\begin{solution}[17.5]
	Subtracting from $E_0(L;\alpha)$ the vacuum energy density and removing the cutoff (taking the limit $\alpha\to0)$, we obtain
	
	$$\Delta E_{\mathrm{ren}}(L)=\lim_{\alpha\to0}\left[E_0\left(L;\alpha\right)-\lim_{L\to\infty}E_0\left(L;\alpha\right)\right]=-\frac\pi{24L}.$$
	
	After we have decided to fix the normalization point and to attribute a “zero” energy to vacuum fluctuations in Minkowski spacetime, there remains no more freedom to renormalize the finite shift of the energy density due to the presence of the plates. 
	
	In the parallel-plate setup (here we consider, more generally, parallel plates in arbitrary dimension, and we denote by $\perp$ the directions along which Neumann boundary conditions are imposed), the allowed local terms are either a bulk cosmological constant term, $\int_{\mathrm{bulk}} d^{d}x$, which implies $E \propto V_\perp L$, or boundary tensions terms, $\int_{\partial} d^{d-1}x,\tau$, which give $E \propto V_\perp$. By contrast, the finite contribution we find, $E \propto L^{-1}$, cannot be removed by any local counterterm, and is therefore entirely physical and observable.
\end{solution}
\begin{tcolorbox}[problem]
\textbf{6.} Generalize to $D$ dimensions with parallel plates separated by $L$, write the mode sum/integral structure and identify how the divergent pieces scale with the transverse volume.
\end{tcolorbox}
\begin{solution}[17.6]
	Consider two parallel plates located at $z=0$ and $z=L$. Based on the analysis in Problem 1, the boundary conditions impose constraints on $k_z$
	\[
	k_z=\frac{n\pi}{L},\quad n=1,2,\ldots
	\]
	However, the transverse directions remain as free plane waves and are unrestricted. By summing over the discrete modes in the $z$-direction and integrating over the phase space of the transverse directions, we obtain the vacuum zero-point energy as:
	\[
	E_0(L)=\frac{1}{2}V_\perp\sum_{n=1}^\infty\int\frac{d^{D-2}k_\perp}{(2\pi)^{D-2}}\sqrt{k_\perp^2+\left(\frac{n\pi}{L}\right)^2}
	\]
	Next, choosing a cutoff, we obtain:
	\begin{align*}
		\frac{E_0(L)}{V_\perp}
		&=\frac12\sum_{n=1}^\infty \int\frac{d^{d}k}{(2\pi)^{d}}
		\sqrt{k^2+\left(\frac{n\pi}{L}\right)^2} \\[4pt]
		&=\frac12\sum_{n=1}^\infty \int\frac{d^{d}k}{(2\pi)^{d}}
		\left[-\frac1{2\sqrt\pi}\int_{1/\Lambda^2}^{\infty}ds\,s^{-3/2}
		\exp\!\left(-s\left[k^2+\left(\frac{n\pi}{L}\right)^2\right]\right)\right] \\[6pt]
		&=-\frac1{4\sqrt\pi}\int_{1/\Lambda^2}^{\infty}ds\,s^{-3/2}
		\sum_{n=1}^\infty \int\frac{d^{d}k}{(2\pi)^{d}}
		e^{-sk^2}\,e^{-s(n\pi/L)^2} \\[6pt]
		&=-\frac1{4\sqrt\pi}\int_{1/\Lambda^2}^{\infty}ds\,s^{-3/2}
		\sum_{n=1}^\infty \frac1{(4\pi s)^{d/2}}\,e^{-s(n\pi/L)^2} \\[6pt]
		&=-\frac1{4\sqrt\pi(4\pi)^{d/2}}
		\int_{1/\Lambda^2}^{\infty}ds\,s^{-(d+3)/2}\sum_{n=1}^\infty e^{-s(n\pi/L)^2} \\[8pt]
		&=-\frac1{4\sqrt\pi(4\pi)^{d/2}}
		\int_{1/\Lambda^2}^{\infty}ds\,s^{-(d+3)/2}
		\left[\frac{L}{2\sqrt{\pi s}}-\frac12+\frac{L}{\sqrt{\pi s}}\sum_{m=1}^{\infty}e^{-m^2L^2/s}\right] \\[10pt]
		&=-\frac{L}{8\pi(4\pi)^{d/2}}\int_{1/\Lambda^2}^{\infty}ds\,s^{-(d+4)/2}
		+\frac{1}{8\sqrt\pi(4\pi)^{d/2}}\int_{1/\Lambda^2}^{\infty}ds\,s^{-(d+3)/2} \\[4pt]
		&\quad-\frac{L}{4\pi(4\pi)^{d/2}}\sum_{m=1}^\infty\int_{1/\Lambda^2}^{\infty}ds\,s^{-(d+4)/2}e^{-m^2L^2/s} \\[10pt]
		&=-\frac{L}{8\pi(4\pi)^{d/2}}\left[\frac{2}{D}\Lambda^{D}\right]
		+\frac{1}{8\sqrt\pi(4\pi)^{d/2}}\left[\frac{2}{D-1}\Lambda^{D-1}\right] \\[4pt]
		&\quad-\frac{L}{4\pi(4\pi)^{d/2}}\sum_{m=1}^\infty\int_{0}^{\infty}ds\,s^{-(d+4)/2}e^{-m^2L^2/s} \\[10pt]
		&=-\frac{L}{D(4\pi)^{D/2}}\Lambda^{D}
		+\frac{1}{2(D-1)(4\pi)^{(D-1)/2}}\Lambda^{D-1} \\[4pt]
		&\quad-\frac{L}{4\pi(4\pi)^{d/2}}\sum_{m=1}^\infty
		\left[\Gamma\!\left(\frac{D}{2}\right)m^{-D}L^{-D}\right] \\[10pt]
		&=-\frac{L}{D(4\pi)^{D/2}}\Lambda^{D}
		+\frac{1}{2(D-1)(4\pi)^{(D-1)/2}}\Lambda^{D-1}
		-\frac{\Gamma(D/2)\,\zeta(D)}{(4\pi)^{D/2}}\frac{1}{L^{D-1}}.
	\end{align*}
	In the above derivation, we used following well-known formula:
	\[
	\begin{gathered}
			\sqrt{A}=-\frac{1}{2\sqrt{\pi}}\int_{1/\Lambda^2}^\infty dss^{-3/2}e^{-sA}\\
			\sum_{n=1}^\infty e^{-s(n\pi/L)^2}=\frac{L}{2\sqrt{\pi s}}-\frac{1}{2}+\frac{L}{\sqrt{\pi s}}\sum_{m=1}^\infty e^{-m^2L^2/s}\\
			\int_0^\infty dss^{-(d+4)/2}e^{-m^2L^2/s}=\Gamma\left(\frac{D}{2}\right)m^{-D}L^{-D}
	\end{gathered}
	\]
	So, \textbf{MY FINAL ANSWER IS}:
	\[
	\boxed{
	E_0(L)=\underbrace{-\frac{V_\perp L}{D(4\pi)^{D/2}}\Lambda^D}_{\text{bulk}}+\underbrace{\frac{V_\perp}{2(D-1)(4\pi)^{(D-1)/2}}\Lambda^{D-1}}_{\text{surface}}\underbrace{-\frac{\Gamma(D/2)\zeta(D)}{(4\pi)^{D/2}}\frac{V_\perp}{L^{D-1}}}_{\text{Casimir}}
	}
	\]
	Based on the analysis in the previous problem, the first two divergent terms can be removed by local counterterms coming from the cosmological constant and the boundary tensions, respectively, while the last term is physically observable. However, in $D=2$ the term associated with the boundary tensions seems to be nonzero, which appears to contradict the conclusion in the previous problem. This is because we adopted different cutoff prescriptions in the two problems. The divergent pieces are naturally scheme-dependent and therefore depend on how the cutoff is implemented, whereas the final finite term is a scheme-independent, physically observable quantity.
\end{solution}
\begin{tcolorbox}[problem]
\textbf{7.} Using dimensional analysis, determine how the finite Casimir energy per transverse area scales with $L$ in general $D$ for a massless field.
\end{tcolorbox}
\begin{solution}[17.7]
From the analysis in the previous problem, we already know that the correct answer should be proportional to $1/L^{D-1}$. Here, we re-derive this result using dimensional analysis. We adopt natural units where $\hbar=c=1$. The dimension of energy is $[E]=[L]^{-1}$. The dimension of the transverse area is $[V_\perp]=[L]^{D-2}$. Therefore, the dimension of energy per unit transverse area is
\begin{equation}
	\left[\frac{E}{V_\perp}\right] = [L]^{-1-(D-2)} = [L]^{-(D-1)}.
\end{equation}
Since the field is massless and characterized by a single length scale $L$, it follows that
\begin{equation}
	\boxed{\frac{E_{0}}{V_{\perp}} \propto \frac{1}{L^{D-1}}}
\end{equation}
\end{solution}

\begin{tcolorbox}[problem]
\textbf{8.} In $D=4$, compute (or quote and justify) the sign of the Casimir force for a massless scalar with Dirichlet boundary conditions between plates.
\end{tcolorbox}
\begin{solution}[17.8]
	From Problem 6, we know the Casimir energy:
	\[
	E_0^{D=4}(L)=-\frac{\Gamma{4/2}\xi(4)}{(4\pi)^2}\frac{V_\perp}{L^3}=-\frac{\pi^2}{1440}\frac{V_\perp}{L^3}
	\]
	So the Casimir force per unit area is given by:
	\[
	\frac{F}{V_\perp}=-\frac{d}{dL}\frac{E_0^{D=4}(L)}{V_\perp}=-\frac{d}{dL}\left(-\frac{\pi^{2}}{1440}L^{-3}\right)=-\frac{\pi^{2}}{480L^{4}}<0
	\]
	This indicates that the Casimir force between the two plates is attractive.
\end{solution}
\begin{tcolorbox}[problem]
\textbf{9.} Replace the scalar with the photon, explain how gauge invariance and physical polarizations imply an effective factor of $D-2$ compared to a scalar (up to boundary condition subtleties).
\end{tcolorbox}

\begin{tcolorbox}[problem]
\textbf{10.} Experimental connection in $D=4$, list the dominant real-world corrections one must control (finite conductivity, surface roughness, temperature, geometry) and explain why Casimir measurements are still a clean test of a QFT vacuum effect rather than a scattering cross section.
\end{tcolorbox}


\end{document}

	\clearpage
	\documentclass[11pt]{article}

\usepackage[margin=1.6cm]{geometry}
\usepackage{amsmath,amssymb,amsthm}
\usepackage[hidelinks]{hyperref}
\usepackage[most]{tcolorbox}
\usepackage[]{cancel}
\usepackage{tikz-feynman}
\tikzfeynmanset{compat=1.1.0}
\newenvironment{solution}[1][]{\begin{proof}[\textbf{Solution}]}{\end{proof}}
% --- Rounded box style for each problem (whole statement) ---
\tcbset{
    problem/.style={
        enhanced,
        boxrule=0.7pt,
        arc=3mm,
        left=3mm,right=3mm,top=2mm,bottom=2mm,
        colback=white,
        colframe=black,
        before skip=10pt,
        after skip=6pt,
    }
}

\begin{document}
\begin{center}
    {\LARGE \textbf{Solution Manual of Problem Set 18} \par}
    \vspace{1em}
    {\large Bufan Zheng\\ \href{mailto:bufan@g.ecc.u-tokyo.ac.jp}{\texttt{bufan@g.ecc.u-tokyo.ac.jp}} \par}
\end{center}


\begin{tcolorbox}[problem]
\textbf{1.} Explain what degeneracy the Lamb shift breaks in the hydrogen spectrum (which states and why they were degenerate in the Dirac theory).
\end{tcolorbox}
\begin{solution}[18.1]
The hydrogen spectrum in Dirac theory can be solved exactly:
\[
E=\frac{mc^2}{\sqrt{1+\frac{\alpha^2}{\left(n-\left|j+\frac12\right|+\sqrt{(j+\frac12)^2-\alpha^2}\right)^2}}}
\]
This shows that the energy levels depend only on the principal quantum number $n$ and the total angular momentum quantum number $j$, and are independent of the orbital angular momentum quantum number $\ell$ and the magnetic quantum number $m_j$.

$2S_{1/2}$ and $2P_{1/2}$ have the same $n$ and $j$, and differ only in the orbital angular momentum quantum number $\ell$. Therefore, in the Dirac theory these two levels should be degenerate. However, the Dirac theory is based on the Coulomb interaction, and in quantum field theory vertex corrections induce a modification to the Coulomb interaction. This correction breaks the degeneracy between these two levels, which constitutes part of the Lamb shift.
\end{solution}
\begin{tcolorbox}[problem]
\textbf{2.} Define the physical observable, energy difference $E(2S_{1/2})-E(2P_{1/2})$ (or equivalent). Explain why it is a non-scattering observable.
\end{tcolorbox}

\begin{tcolorbox}[problem]
\textbf{3.} Identify the two main QED contributions at leading order, electron self-energy in the Coulomb field and vacuum polarization. State which dominates and why.
\end{tcolorbox}

\begin{tcolorbox}[problem]
\textbf{4.} Explain the origin of the ``Bethe logarithm'' qualitatively, which range of virtual photon energies contributes and why a log appears.
\end{tcolorbox}

\begin{tcolorbox}[problem]
\textbf{5.} Show (by dimensional analysis) why bound-state radiative corrections naturally involve $\ln(m/\alpha m)=\ln(1/\alpha)$.
\end{tcolorbox}

\begin{tcolorbox}[problem]
\textbf{6.} In dimensional regularization, explain how UV divergences in the self-energy are renormalized in the same way as in scattering, yet the finite remainder depends on the bound-state environment.
\end{tcolorbox}

\begin{tcolorbox}[problem]
\textbf{7.} Explain why an effective field theory viewpoint (NRQED) is helpful, separate ``hard'' ($\sim m$), ``soft'' ($\sim \alpha m$), and ``ultrasoft'' ($\sim \alpha^2 m$) scales.
\end{tcolorbox}

\begin{tcolorbox}[problem]
\textbf{8.} Compute a simple toy version, for a contact perturbation $\delta V(r)=c\,\delta^{(3)}(r)$, show it shifts only S-wave states and estimate the scaling with principal quantum number $n$.
\end{tcolorbox}
\begin{solution}[18.8]
	According to time-independent perturbation theory, the leading-order correction to the hydrogen energy levels is given by   
	\[
	\Delta E_{n\ell m}^{(1)}=\langle n\ell m|\delta V|n\ell m\rangle=c\int d^3r\left|\psi_{n\ell m}(\mathbf{r})\right|^2\delta^{(3)}(\mathbf{r})=c\left|\psi_{n\ell m}(0)\right|^2.
	\]
	The unperturbed hydrogen wavefunction is
	\[
	\psi_{n\ell m}=\sqrt{\left(\frac{2}{na}\right)^{3}\frac{(n-\ell-1)!}{2n\left(n+\ell\right)!}}e^{-r/na}\left(\frac{2r}{na}\right)^{\ell}\left[L_{n-\ell-1}^{2\ell+1}(2r/na)\right]Y_{\ell}^{m}(\theta,\phi)
	\]
	From this we obtain
	\[
	|\psi_{n00}(0)|^2=\frac{1}{\pi a_0^3}\frac{1}{n^3}
	\]
	Substituting this into the expression for the energy shift derived above, we find 
	\[
	\boxed{\Delta E_{nS}^{(1)}=c\left|\psi_{n00}(0)\right|^2=\frac{c}{\pi a_0^3}\frac{1}{n^3}\propto n^{-3}.}
	\]
\end{solution}
\begin{tcolorbox}[problem]
\textbf{9.} Discuss dimension dependence, if you formally considered ``hydrogen'' in $D\neq 4$, which parts of the Lamb-shift analysis would fail or change qualitatively? (Potential law, density of states, IR/UV behavior.)
\end{tcolorbox}

\begin{tcolorbox}[problem]
\textbf{10.} Conceptual summary, explain in one paragraph why the Lamb shift is a triumph of QFT renormalization (local counterterms) yet produces a finite, physical, environment-dependent energy shift that cannot be removed by redefining parameters.
\end{tcolorbox}


\end{document}

	
\end{document}